\IfFileExists{stacks-project.cls}{%
\documentclass{stacks-project}
}{%
\documentclass{amsart}
}
\usepackage{amssymb}

% For dealing with references we use the comment environment
\usepackage{verbatim}
\usepackage{fancyvrb}
\usepackage{pmboxdraw}
\newenvironment{reference}{\comment}{\endcomment}
%\newenvironment{reference}{}{}
\newenvironment{slogan}{\comment}{\endcomment}
\newenvironment{history}{\comment}{\endcomment}

% For source figures retained as assets
\usepackage{graphicx}

% For commutative diagrams we use Xy-pic
\usepackage[all]{xy}

% We use 2cell for 2-commutative diagrams.
\xyoption{2cell}
\UseAllTwocells

% We use multicol for the list of chapters between chapters
\usepackage{multicol}

% This is generally recommended for better output
\usepackage{lmodern}
\usepackage[T1]{fontenc}
\usepackage{textalpha}

% For cross-file-references
% We prefer the snapshot tmp/xref/#2.aux,
% which latexmkrc replaces atomically, so
% that documents compiled simultaneously
% never read a half-written .aux file.
% The last argument keeps links to #2.pdf.
\usepackage{xr-hyper}
\newcommand{\maybeexternaldocument}[2]{%
  \IfFileExists{#2.aux}{%
    \externaldocument[#1]{#2}%
  }{%
    \PackageWarning{stack}{Missing #2.aux; external references with prefix #1 unavailable}%
  }%
}
\let\externaldocumentorig\externaldocument
\renewcommand{\externaldocument}[2][]{%
  \IfFileExists{tmp/xref/#2.aux}{%
    \externaldocumentorig[#1]{tmp/xref/#2}[#2.pdf]%
  }{%
  \IfFileExists{#2.aux}{%
    \externaldocumentorig[#1]{#2}%
  }{%
    \PackageWarning{stack}{Missing #2.aux; external references with prefix #1 unavailable}%
  }%
  }%
}

% Package for hypertext links:
\usepackage{hyperref}

% For any local file, say "hello.tex" you want to link to please
% use \externaldocument[hello-]{hello}
\externaldocument[introduction-]{introduction}
\externaldocument[conventions-]{conventions}
\externaldocument[sets-]{sets}
\externaldocument[categories-]{categories}
\externaldocument[topology-]{topology}
\externaldocument[sheaves-]{sheaves}
\externaldocument[sites-]{sites}
\externaldocument[stacks-]{stacks}
\externaldocument[fields-]{fields}
\externaldocument[algebra-]{algebra}
\externaldocument[brauer-]{brauer}
\externaldocument[homology-]{homology}
\externaldocument[derived-]{derived}
\externaldocument[simplicial-]{simplicial}
\externaldocument[more-algebra-]{more-algebra}
\externaldocument[smoothing-]{smoothing}
\externaldocument[modules-]{modules}
\externaldocument[sites-modules-]{sites-modules}
\externaldocument[injectives-]{injectives}
\externaldocument[cohomology-]{cohomology}
\externaldocument[sites-cohomology-]{sites-cohomology}
\externaldocument[dga-]{dga}
\externaldocument[dpa-]{dpa}
\externaldocument[sdga-]{sdga}
\externaldocument[hypercovering-]{hypercovering}
\externaldocument[schemes-]{schemes}
\externaldocument[constructions-]{constructions}
\externaldocument[properties-]{properties}
\externaldocument[morphisms-]{morphisms}
\externaldocument[coherent-]{coherent}
\externaldocument[divisors-]{divisors}
\externaldocument[limits-]{limits}
\externaldocument[varieties-]{varieties}
\externaldocument[topologies-]{topologies}
\externaldocument[descent-]{descent}
\externaldocument[perfect-]{perfect}
\externaldocument[more-morphisms-]{more-morphisms}
\externaldocument[flat-]{flat}
\externaldocument[groupoids-]{groupoids}
\externaldocument[more-groupoids-]{more-groupoids}
\externaldocument[etale-]{etale}
\externaldocument[chow-]{chow}
\externaldocument[intersection-]{intersection}
\externaldocument[pic-]{pic}
\externaldocument[weil-]{weil}
\externaldocument[adequate-]{adequate}
\externaldocument[dualizing-]{dualizing}
\externaldocument[duality-]{duality}
\externaldocument[discriminant-]{discriminant}
\externaldocument[derham-]{derham}
\externaldocument[local-cohomology-]{local-cohomology}
\externaldocument[algebraization-]{algebraization}
\externaldocument[curves-]{curves}
\externaldocument[resolve-]{resolve}
\externaldocument[models-]{models}
\externaldocument[functors-]{functors}
\externaldocument[equiv-]{equiv}
\externaldocument[pione-]{pione}
\externaldocument[etale-cohomology-]{etale-cohomology}
\externaldocument[proetale-]{proetale}
\externaldocument[relative-cycles-]{relative-cycles}
\externaldocument[more-etale-]{more-etale}
\externaldocument[trace-]{trace}
\externaldocument[crystalline-]{crystalline}
\externaldocument[spaces-]{spaces}
\externaldocument[spaces-properties-]{spaces-properties}
\externaldocument[spaces-morphisms-]{spaces-morphisms}
\externaldocument[decent-spaces-]{decent-spaces}
\externaldocument[spaces-cohomology-]{spaces-cohomology}
\externaldocument[spaces-limits-]{spaces-limits}
\externaldocument[spaces-divisors-]{spaces-divisors}
\externaldocument[spaces-over-fields-]{spaces-over-fields}
\externaldocument[spaces-topologies-]{spaces-topologies}
\externaldocument[spaces-descent-]{spaces-descent}
\externaldocument[spaces-perfect-]{spaces-perfect}
\externaldocument[spaces-more-morphisms-]{spaces-more-morphisms}
\externaldocument[spaces-flat-]{spaces-flat}
\externaldocument[spaces-groupoids-]{spaces-groupoids}
\externaldocument[spaces-more-groupoids-]{spaces-more-groupoids}
\externaldocument[bootstrap-]{bootstrap}
\externaldocument[spaces-pushouts-]{spaces-pushouts}
\externaldocument[spaces-chow-]{spaces-chow}
\externaldocument[groupoids-quotients-]{groupoids-quotients}
\externaldocument[spaces-more-cohomology-]{spaces-more-cohomology}
\externaldocument[spaces-simplicial-]{spaces-simplicial}
\externaldocument[spaces-duality-]{spaces-duality}
\externaldocument[formal-spaces-]{formal-spaces}
\externaldocument[restricted-]{restricted}
\externaldocument[spaces-resolve-]{spaces-resolve}
\externaldocument[formal-defos-]{formal-defos}
\externaldocument[defos-]{defos}
\externaldocument[cotangent-]{cotangent}
\externaldocument[examples-defos-]{examples-defos}
\externaldocument[algebraic-]{algebraic}
\externaldocument[examples-stacks-]{examples-stacks}
\externaldocument[stacks-sheaves-]{stacks-sheaves}
\externaldocument[criteria-]{criteria}
\externaldocument[artin-]{artin}
\externaldocument[quot-]{quot}
\externaldocument[stacks-properties-]{stacks-properties}
\externaldocument[stacks-morphisms-]{stacks-morphisms}
\externaldocument[stacks-limits-]{stacks-limits}
\externaldocument[stacks-cohomology-]{stacks-cohomology}
\externaldocument[stacks-perfect-]{stacks-perfect}
\externaldocument[stacks-introduction-]{stacks-introduction}
\externaldocument[stacks-more-morphisms-]{stacks-more-morphisms}
\externaldocument[stacks-geometry-]{stacks-geometry}
\externaldocument[moduli-]{moduli}
\externaldocument[moduli-curves-]{moduli-curves}
\externaldocument[examples-]{examples}
\externaldocument[exercises-]{exercises}
\externaldocument[guide-]{guide}
\externaldocument[desirables-]{desirables}
\externaldocument[coding-]{coding}
\externaldocument[obsolete-]{obsolete}
\externaldocument[fdl-]{fdl}
\externaldocument[index-]{index}

% Heap Project documents
\externaldocument[k3-]{k3}
\externaldocument[torelli-]{torelli}
\externaldocument[linear-algebra-]{linear-algebra}
\externaldocument[groups-]{groups}
\externaldocument[complex-analysis-]{complex-analysis}
\externaldocument[functional-analysis-]{functional-analysis}
\externaldocument[categories-]{categories}
\externaldocument[commutative-algebra-]{commutative-algebra}
\externaldocument[ringed-spaces-]{ringed-spaces}
\externaldocument[homological-algebra-]{homological-algebra}
\externaldocument[lie-algebras-]{lie-algebras}
\externaldocument[representation-theory-]{representation-theory}
\externaldocument[smooth-manifolds-]{smooth-manifolds}
\externaldocument[differential-topology-]{differential-topology}
\externaldocument[algebraic-topology-]{algebraic-topology}
\externaldocument[differential-geometry-]{differential-geometry}
\externaldocument[algebraic-geometry-]{algebraic-geometry}
\externaldocument[symplectic-geometry-]{symplectic-geometry}
\externaldocument[complex-geometry-]{complex-geometry}
\externaldocument[hodge-theory-]{hodge-theory}
\externaldocument[%
variations-of-hodge-structures-]{%
variations-of-hodge-structures}
\externaldocument[%
homological-mirror-symmetry-]{%
homological-mirror-symmetry}
\maybeexternaldocument{deformations-}{deformations}
\externaldocument[vertex-algebras-]{vertex-algebras}
\externaldocument[springer-]{springer}
\externaldocument[infty-categories-]{infty-categories}
\externaldocument[seminars-differential-geometry-]{%
seminars-differential-geometry}
\externaldocument[seminars-complex-geometry-]{%
seminars-complex-geometry}
\externaldocument[seminars-algebraic-geometry-]{%
seminars-algebraic-geometry}
\externaldocument[seminars-analysis-pde-]{%
seminars-analysis-pde}
\externaldocument[surfaces-]{surfaces}
\externaldocument[birational-maps-commutative-algebra-]{birational-maps-commutative-algebra}
\externaldocument[cremona-transformations-]{cremona-transformations}
\externaldocument[derived-categories-of-sheaves-]{derived-categories-of-sheaves}
\externaldocument[geometric-invariant-theory-]{geometric-invariant-theory}
\externaldocument[geometric-stability-conditions-and-group-actions-]{geometric-stability-conditions-and-group-actions}
\externaldocument[moduli-spaces-of-sheaves-]{moduli-spaces-of-sheaves}
\externaldocument[bridgeland-stability-general-theory-]{bridgeland-stability-general-theory}
\externaldocument[hilbert-schemes-]{hilbert-schemes}
\externaldocument[noncommutative-abelian-surfaces-and-kummer-type-hyperkahler-varieties-]{noncommutative-abelian-surfaces-and-kummer-type-hyperkahler-varieties}
\externaldocument[mumford-tate-groups-in-hodge-theory-]{mumford-tate-groups-in-hodge-theory}
\externaldocument[hodge-birational-atoms-2026-]{hodge-birational-atoms-2026}
\externaldocument[xxii-egd-2026-]{%
xxii-egd-2026}
\externaldocument[pde-]{pde}
\externaldocument[probability-]{probability}

\externaldocument[linguistics-]{linguistics}

\externaldocument[%
16th-alga-meeting-2026-]{%
16th-alga-meeting-2026}

\externaldocument[%
iv-jornadas-lefschetz-2026-]{%
iv-jornadas-lefschetz-2026}

\externaldocument[reading-the-masters-]{%
reading-the-masters}

\externaldocument[real-analysis-]{%
real-analysis}

\externaldocument[optimal-transport-]{%
optimal-transport}

\externaldocument[history-]{history}

\externaldocument[seiberg-witten-]{seiberg-witten}

\externaldocument[history-philosophy-biology-]{history-philosophy-biology}

\externaldocument[classical-mechanics-]{classical-mechanics}
\externaldocument[electromagnetism-]{electromagnetism}
\externaldocument[waves-and-optics-]{waves-and-optics}
\externaldocument[wif-]{waves-in-fluids-iv}
\externaldocument[thermal-physics-]{thermal-physics}
\externaldocument[quantum-mechanics-]{quantum-mechanics}
\externaldocument[special-relativity-]{special-relativity}
\externaldocument[physics-]{physics}

\externaldocument[chemistry-]{chemistry}

\externaldocument[molecular-cell-biology-]{molecular-cell-biology}

\externaldocument[microbiology-]{microbiology}

\externaldocument[protistology-]{protistology}

\externaldocument[mycology-]{mycology}

\externaldocument[botany-]{botany}

\externaldocument[zoology-]{zoology}

\externaldocument[animal-histology-]{animal-histology}

\externaldocument[animal-physiology-]{animal-physiology}

\externaldocument[systematics-biogeography-]{systematics-biogeography}

\externaldocument[ecology-natural-resources-]{ecology-natural-resources}

\externaldocument[earth-sciences-]{earth-sciences}

\externaldocument[biostatistics-]{biostatistics}

\externaldocument[mathematical-systems-biology-]{mathematical-systems-biology}

\externaldocument[cimpa-el-salvador-2024-]{cimpa-el-salvador-2024}

\externaldocument[anthropology-]{anthropology}

\externaldocument[yang-mills-theory-]{yang-mills-theory}

\externaldocument[seminars-algebra-]{seminars-algebra}

\externaldocument[seminars-physics-]{seminars-physics}

\externaldocument[seminars-computing-]{seminars-computing}

\externaldocument[oktobermat-2026-]{oktobermat-2026}

\externaldocument[saet-2026-]{saet-2026}

\externaldocument[seminars-symplectic-geometry-]{seminars-symplectic-geometry}

\externaldocument[seminars-dynamics-]{seminars-dynamics}

\externaldocument[oktobermat-2024-]{oktobermat-2024}

\externaldocument[quantum-groups-]{quantum-groups}

\externaldocument[legal-2027-]{legal-2027}

\externaldocument[derived-categories-]{derived-categories}

% Theorem environments.
%
\theoremstyle{plain}
\newtheorem{theorem}[subsection]{Theorem}
\newtheorem{proposition}[subsection]{Proposition}
\newtheorem{lemma}[subsection]{Lemma}
\newtheorem{corollary}[subsection]{Corollary}

\theoremstyle{definition}
\newtheorem{definition}[subsection]{Definition}
\newtheorem{example}[subsection]{Example}
\newtheorem{exercise}[subsection]{Exercise}
\newtheorem{situation}[subsection]{Situation}

\theoremstyle{remark}
\newtheorem{remark}[subsection]{Remark}
\newtheorem{remarks}[subsection]{Remarks}

\numberwithin{equation}{subsection}

% Macros
%
\def\lim{\mathop{\mathrm{lim}}\nolimits}
\def\colim{\mathop{\mathrm{colim}}\nolimits}
\def\Spec{\mathop{\mathrm{Spec}}}
\def\Hom{\mathop{\mathrm{Hom}}\nolimits}
\def\Ext{\mathop{\mathrm{Ext}}\nolimits}
\def\SheafHom{\mathop{\mathcal{H}\!\mathit{om}}\nolimits}
\def\SheafExt{\mathop{\mathcal{E}\!\mathit{xt}}\nolimits}
\def\Sch{\mathit{Sch}}
\def\Mor{\mathop{\mathrm{Mor}}\nolimits}
\def\Ob{\mathop{\mathrm{Ob}}\nolimits}
\def\Sh{\mathop{\mathit{Sh}}\nolimits}
\def\NL{\mathop{N\!L}\nolimits}
\def\CH{\mathop{\mathrm{CH}}\nolimits}
\def\proetale{{pro\text{-}\acute{e}tale}}
\def\etale{{\acute{e}tale}}
\def\QCoh{\mathit{QCoh}}
\def\Ker{\mathop{\mathrm{Ker}}}
\def\Im{\mathop{\mathrm{Im}}}
\def\Coker{\mathop{\mathrm{Coker}}}
\def\Coim{\mathop{\mathrm{Coim}}}

% Boxtimes
%
\DeclareMathSymbol{\boxtimes}{\mathbin}{AMSa}{"02}

%
% Macros for moduli stacks/spaces
%
\def\QCohstack{\mathcal{QC}\!\mathit{oh}}
\def\Cohstack{\mathcal{C}\!\mathit{oh}}
\def\Spacesstack{\mathcal{S}\!\mathit{paces}}
\def\Quotfunctor{\mathrm{Quot}}
\def\Hilbfunctor{\mathrm{Hilb}}
\def\Curvesstack{\mathcal{C}\!\mathit{urves}}
\def\Polarizedstack{\mathcal{P}\!\mathit{olarized}}
\def\Complexesstack{\mathcal{C}\!\mathit{omplexes}}
% \Pic is the operator that assigns to X its picard group, usage \Pic(X)
% \Picardstack_{X/B} denotes the Picard stack of X over B
% \Picardfunctor_{X/B} denotes the Picard functor of X over B
\def\Pic{\mathop{\mathrm{Pic}}\nolimits}
\def\Picardstack{\mathcal{P}\!\mathit{ic}}
\def\Picardfunctor{\mathrm{Pic}}
\def\Deformationcategory{\mathcal{D}\!\mathit{ef}}


\begin{document}

\title{Topological Groups}
\maketitle

\phantomsection
\label{section-phantom}

\tableofcontents

\section{Overview}
\label{section-overview}

\noindent
These are my notes from a course on
topological transformation groups. They
were written in Spanish under the title
{\it Grupos topol\'ogicos de
transformaciones}; the original is kept
in {\tt standalone/} as
{\tt grupos-topologicos-transformaciones}.
The source does not name the course, the
lecturer or the semester. This chapter is
a translation that keeps the content as
it was: statements without proof stay
without proof, unfinished proofs stay
unfinished, and questions stay open.
Where the original is wrong, incomplete
or unclear, a sentence starting with
{\it Note on the original} says so.
Statements marked ``for presentation''
were topics for student presentations.

\medskip\noindent
{\bf Big picture.} A topological group is
a group with a topology for which the
group operations are continuous. It acts
on a space $X$ by homeomorphisms, and
$X$ is then called a $G$-space. The
notes go through the following steps.
\begin{enumerate}
\item Topological groups are homogeneous
and regular
(Sections
\ref{section-definition-examples} and
\ref{section-topological-properties}).
\item Orbits, stabilizers and the orbit
space $X/G$; for compact $G$ the orbit
space inherits the good properties of
$X$, including metrizability
(Section \ref{section-group-actions}).
\item The twisted product
$G\times_HX$ builds a $G$-space out of
an $H$-space
(Section \ref{section-twisted-product}).
\item A slice is a subset whose
saturation is a twisted product; slices
give the local structure of a $G$-space
(Section \ref{section-slices}).
\item Extension of equivariant maps:
averaging with the Haar measure, the
theorem of Dugundji and its equivariant
version through $G$-canonical covers
(Sections \ref{section-extension},
\ref{section-dugundji} and
\ref{section-canonical-cover}).
\item Proper actions of locally compact
groups
(Section \ref{section-proper-actions}).
\end{enumerate}

\section{Definition and examples}
\label{section-definition-examples}

\begin{definition}
\label{definition-topological-group}
A {\it topological group} is a triple
$(G,\tau,\cdot)$ where $\tau$ is a
topology and $\cdot$ is a group
operation such that the multiplication
$$
G\times G\to G,
\qquad
(x,y)\mapsto x\cdot y
$$
and the symmetry (inversion)
$$
G\to G,
\qquad
x\mapsto -x
$$
are continuous.
\end{definition}

\noindent
{\it Note on the original.} The inversion
is written additively there; in
multiplicative notation it is
$x\mapsto x^{-1}$.

\begin{example}
\label{example-first-topological-groups}
\begin{enumerate}
\item $(\mathbb{R},+)$.
\item $(\mathbb{R}^n,+)$. Exercise: prove
with $\varepsilon$ and $\delta$ that the
sum is continuous. Solution: let
$\varepsilon>0$ and let $(x,y)$ and
$(x',y')$ be points of
$\mathbb{R}^n\times\mathbb{R}^n$.
Suppose that both lie in an open square
of side less than $\varepsilon/2$. Then
$$
\begin{aligned}
|x+y-(x'+y')|
&=|x-x'+y-y'| \\
&\leq|x-x'|+|y-y'| \\
&<\varepsilon.
\end{aligned}
$$
\item If $X$ is a normed linear space,
then $(X,+)$ is a topological group.
\item $(\mathbb{R}_+,\cdot)$ where
$$
\mathbb{R}_+
=\{x\in\mathbb{R}:x>0\}.
$$
In this case the space is not a vector
space (it is isomorphic to one, but
that is not obvious for now).
\item
$(\mathbb{R}\setminus\{0\},\cdot)$,
which is not connected.
\item $(\operatorname{GL}(n),\cdot)$,
where $\operatorname{GL}(n)$ is the set
of invertible matrices $A=(a_{ij})$ with
real entries $a_{ij}$ and the operation
is the usual product of matrices.

Recall that a matrix is invertible if
and only if its determinant is not zero.
Since the determinant is multiplicative,
the product is well defined, so this is
indeed a group. The topology is the one
induced by the metric of
$\mathbb{R}^{n^2}$.

Now we check that the group operations
are continuous. For the product: the
product $AB=C$ of two matrices, with
$C=(c_{ij})$, is given by
$$
c_{ij}=\sum_{k=1}^na_{ik}b_{kj},
$$
which is a continuous function.

For the inversion, suppose that
$A^{-1}=(c_{ij})$, so that
$$
c_{ij}
=\frac{(-1)^{i+j}a_{ij}}{\det A_{ij}}
$$
where $A_{ij}$ is the matrix of minors
obtained by removing row $i$ and
column $j$. Now consider
$$
\xymatrix{
\operatorname{GL}(n,\mathbb{R})
  \ar@{^{(}->}[d] \ar[r]^-{\det}
  & \mathbb{R} \\
M(n,\mathbb{R})
  \approx\mathbb{R}^{n^2}
  \ar[ur]_{\det}
}
$$
and since the determinant is a
polynomial, continuity follows.

{\it Note on the original.} The formula
for $c_{ij}$ is not the formula for the
inverse matrix. The correct one is
$$
c_{ij}
=\frac{(-1)^{i+j}\det A_{ji}}{\det A},
$$
a quotient of polynomials in the entries
with nonvanishing denominator, which is
what the argument needs.
\end{enumerate}
\end{example}

\begin{definition}
\label{definition-lie-group}
A topological group $G$ is called a
{\it Lie group} if
\begin{enumerate}
\item $G$ is a differentiable manifold,
\item the multiplication and the
inversion are differentiable.
\end{enumerate}
\end{definition}

\begin{remark}
\label{remark-lie-group-examples}
All the examples so far are Lie groups.
An example that is not a Lie group is
$C([0,1])$, the continuous functions on
$[0,1]$, since it is not locally
homeomorphic to $\mathbb{R}^n$ with $n$
finite.
\end{remark}

\begin{definition}
\label{definition-topological-subgroup}
Let $G$ be a topological group and
$H<G$ an algebraic subgroup. With the
induced topology $H$ is a topological
group (since the restriction of a
continuous function is continuous). It
is called a {\it topological subgroup}
of $G$.
\end{definition}

\begin{example}[For presentation]
\label{example-matrix-groups}
\begin{enumerate}
\item $\operatorname{GL}(n,\mathbb{R})$,
since it is an open subset of
$\mathbb{R}^{n^2}$ (it is the preimage
of the open set
$\mathbb{R}\setminus\{0\}$ under the
continuous function $\det$).
\begin{enumerate}
\item $\operatorname{GL}(1)
=\mathbb{R}\setminus\{0\}$.
\item The matrices of
$\operatorname{GL}(n)$ with positive
determinant and with negative
determinant are topological subgroups.
\end{enumerate}
\item The subgroup of orthogonal
matrices
$\operatorname{O}(n,\mathbb{R})
<\operatorname{GL}(n,\mathbb{R})$,
that is, the set of
$A\in\operatorname{GL}(n,\mathbb{R})$
such that
$$
A^{-1}=A^t.
$$
It is also a Lie group which, unlike
$\operatorname{GL}(n,\mathbb{R})$, is
compact.
\begin{enumerate}
\item $\operatorname{O}(1)
=\{-1,1\}=\mathbb{Z}_2$.
\item $\operatorname{SO}(1)=\{1\}$.
\end{enumerate}
\item The special linear group
$\operatorname{SL}(n,\mathbb{R})
<\operatorname{GL}(n,\mathbb{R})$,
that is, the set of
$A\in\operatorname{GL}(n,\mathbb{R})$
such that
$$
\det A=1,
$$
which is not compact either.
\item The group
$$
\operatorname{SO}(n,\mathbb{R})
=\operatorname{O}(n,\mathbb{R})
\cap\operatorname{SL}(n,\mathbb{R}).
$$
According to the GAP program, the
subgroup generated by the identity has
index 2; is this possible? How do we
prove that the two cosets are
disconnected from each other?
\item $\mathbb{C}^*
=\mathbb{C}\setminus\{0\}$, which has
as a subgroup
$$
\mathbb{S}^1
=\{z\in\mathbb{C}:|z|=1\}.
$$
\item $\mathbb{S}^3\subset\mathbb{C}^2$
is also a group, but $\mathbb{S}^2$ is
not. This fact is related to the fact
that $\mathbb{R}^3$ cannot have a field
structure, whereas
$\mathbb{R}^2\approx\mathbb{C}$ and
$\mathbb{R}$ can. For its part,
$\mathbb{R}^4$ seen as the quaternions
is a {\it skew field}, which satisfies
all the field axioms except the
existence of multiplicative inverses.
\end{enumerate}
\end{example}

\noindent
{\it Note on the original.} The matrices
with negative determinant do not form a
subgroup: they do not contain the
identity. In the item on
$\operatorname{SO}(n,\mathbb{R})$ the
question is left as it was asked; the
subgroup of index 2 in
$\operatorname{O}(n,\mathbb{R})$ is
$\operatorname{SO}(n,\mathbb{R})$, the
connected component of the identity. The
quaternions do have multiplicative
inverses; the field axiom that fails is
the commutativity of the product.

\section{Topological properties of
topological groups}
\label{section-topological-properties}

\begin{remark}
\label{remark-translations-open}
As we will see with $G$-spaces, for each
element $g\in G$ we have a
homeomorphism $\tilde{g}$ that sends
each point to its product with $g$. For
any open set $U\subseteq G$, the sets
$g\cdot U$ and $U\cdot g$ are open, and
the same holds for a closed set. The
product of an open set $U$ with any set
$W$, defined pointwise, is open, since
$$
W\cdot U=\bigcup_{w\in W}w\cdot U.
$$
\end{remark}

\begin{definition}
\label{definition-neighborhood}
A subset $V\subseteq X$ of a topological
space is a {\it neighborhood} of
$x\in X$ if $x$ lies in the interior of
$V$. The set of neighborhoods of the
identity $e$ of a topological group $G$
is denoted by $N(e)$.
\end{definition}

\begin{definition}
\label{definition-homogeneous}
A topological space $X$ is
{\it homogeneous} if for any two points
$x,y\in X$ there is a homeomorphism
$f:X\to X$ such that $f(x)=y$.
\end{definition}

\begin{remark}
\label{remark-groups-homogeneous}
Every topological group is homogeneous.
\end{remark}

\begin{proposition}
\label{proposition-local-base-translate}
Let $B$ be a local base at $e$. Then
$$
B_x=\{x\cdot V:V\in B\}
$$
is a local base at any $x\in G$.
\end{proposition}

\begin{proof}
Let $V\in\tau$ with $x\in V$. Then
$e\in x^{-1}V$. Since $B$ is a base,
there is $W\in B$ such that
$$
e\in W\subseteq x^{-1}\cdot V.
$$
Then
$$
x\in x\cdot W
\subseteq x\cdot x^{-1}V=V,
$$
and $x\cdot W\in B_x$.
\end{proof}

\begin{definition}
\label{definition-symmetric-neighborhood}
A set $V\in N(e)$ is a
{\it symmetric neighborhood} if
$V=V^{-1}$.
\end{definition}

\begin{proposition}
\label{proposition-symmetric-neighborhood}
Let $V\in N(e)\cap\tau_G$. Then there is
$W\in N(e)\cap\tau_G$ such that
$$
W=W^{-1}\subset V.
$$
\end{proposition}

\begin{proof}
Let $e\in V\in\tau_G$. Since the
inversion is continuous, there is
$W'\in N(e)\cap\tau_G$ such that
$(W')^{-1}\subset V$. We propose
$W'\cap(W')^{-1}$.
\end{proof}

\noindent
{\it Note on the original.} For
$W'\cap(W')^{-1}\subset V$ one should
also ask $W'\subset V$, which can be
arranged by intersecting $W'$ with $V$.

\begin{remark}
\label{remark-square-root-neighborhood}
If $V\in N(e)\cap\tau$, there is
$W\in N(e)\cap\tau$ such that $W^2=V$.
\end{remark}

\noindent
{\it Note on the original.} The equality
should be the inclusion $W^2\subset V$,
which is the first item of the next
theorem.

\begin{theorem}
\label{theorem-neighborhoods-of-identity}
Let $G$ be a topological group. Then
\begin{enumerate}
\item for every $U\in N(e)\cap\tau$
there is $V\in N(e)\cap\tau$ such that
$V^2\subset U$;
\item for every $U\in N(e)\cap\tau$
there is $V\in N(e)\cap\tau$ such that
$V^2\subset U$ and $V^{-1}=V$;
\item for every $U\in N(e)\cap\tau$
there is $V\in N(e)\cap\tau$ which is
{\it regular}, that is, such that
$V\subset\bar{V}\subset U$.
\end{enumerate}
\end{theorem}

\begin{proof}
(1) Let $U\in N(e)\cap\tau$. Since the
product is continuous and
$e\cdot e=e\in U$, there are two open
sets $V_1,V_2\in N(e)$ such that
$V_1\cdot V_2\subset U$. Define
$V=V_1\cap V_2$. Then
$V\in N(e)\cap\tau$ and
$$
V\cdot V
=(V_1\cap V_2)\cdot(V_1\cap V_2)
\subset V_1\cdot V_2\subset U.
$$

(2) Let $U\in N(e)\cap\tau$. Then there
is $V_1\in N(e)\cap\tau$ with
$V_1^2\subset U$. By
Proposition
\ref{proposition-symmetric-neighborhood}
there is $V\in N(e)\cap\tau$ such that
$V^{-1}=V$ and $V\subset V_1$. Then
$$
V\cdot V\subset V_1\cdot V_1
=V_1^2\subset U.
$$

(3) Let $U\in N(e)\cap\tau$. By the
previous item there is
$V\in N(e)\cap\tau$ such that
$V^2\subset U$ and $V^{-1}=V$. We claim
that $\bar{V}\subset U$. Let
$x\in\bar{V}$. The set $x\cdot V$ is a
neighborhood of $x$, so
$$
x\cdot V\cap V\neq\varnothing.
$$
Thus there are $v_1,v_2\in V$ such that
$x\cdot v_1=v_2$. Then
$$
x=v_2\cdot v_1^{-1}
\in V\cdot V=V^2\subset U,
$$
that is, $x\in U$.
\end{proof}

\noindent
{\it Note on the original.} In item (3)
the original writes ``$x\in x\cdot V$,
so $x\in x\cdot V\cap V$'' and uses
capital letters $V_1,V_2$ for the two
elements of $V$. The version above is
the intended argument.

\begin{definition}
\label{definition-separation-axioms}
The topological space $X$ is
\begin{enumerate}
\item {\it $T_0$} if for any distinct
$x,y\in X$ there is $U\in\tau_X$ such
that $|U\cap\{x,y\}|=1$;
\item {\it $T_1$} if for any distinct
$x,y\in X$ there are open sets $U,V$
such that $x\in U\setminus V$ and
$y\in V\setminus U$;
\item {\it $T_2$} if for any two points
$x,y\in X$ there are open sets $U,V$
such that $x\in U$, $y\in V$ and
$U\cap V\neq\varnothing$;
\item {\it $T_3$} if $X$ is $T_1$ and
moreover for any point $x\in X$ and any
closed set $A\subset X$ there are open
sets $U,V$ such that $x\in U$,
$A\subset V$ and
$U\cap V\neq\varnothing$.
\end{enumerate}
\end{definition}

\noindent
{\it Note on the original.} In $T_2$ and
$T_3$ the condition should be
$U\cap V=\varnothing$, the two points in
$T_2$ should be distinct, and in $T_3$
the point $x$ should not belong to $A$.
The condition of $T_3$ is garbled in the
original; it is reconstructed here.

\begin{remark}
\label{remark-separation-axioms}
\begin{enumerate}
\item $X$ is $T_1$ if and only if
$\{x\}$ is closed for every point
$x\in X$.
\item If $X$ is $T_1$ and regular, then
it is $T_3$.
\item There are implications
$$
T_3\Rightarrow T_2
\Rightarrow T_1\Rightarrow T_0.
$$
\end{enumerate}
\end{remark}

\begin{theorem}
\label{theorem-t0-implies-t3}
Every $T_0$ topological group is $T_3$.
\end{theorem}

\begin{proof}
Since any topological group satisfies
condition (3) of
Theorem
\ref{theorem-neighborhoods-of-identity},
we only have to show that it is $T_1$,
that is, that the one-point sets are
closed. Since $G$ is homogeneous, it is
enough to prove that $\{e\}$ is closed.

We show that if $g\in G\setminus\{e\}$
then there is $V\in\tau_G$ such that
$$
g\in V\subset G\setminus\{e\}.
$$
Since $g\neq e$ and $G$ is $T_0$, there
is an open set $U$ such that
$|U\cap\{g,e\}|=1$. There are two cases.
\begin{enumerate}
\item If $g\in U$, then $e\notin U$, so
$g\in U\subset G\setminus\{e\}$.
\item If $e\in U$, then there is
$V\in N(e)\cap\tau_G$ with
$\bar{V}\subset U$. Clearly
$g\notin\bar{V}$, so
$g\in G\setminus\bar{V}\in\tau_G$ and
$e\notin G\setminus\bar{V}$. Then
$$
g\in G\setminus\bar{V}
\subset G\setminus\{e\}
$$
and $G\setminus\{e\}\in\tau_G$.
\end{enumerate}
\end{proof}

\begin{remark}
\label{remark-translations-no-fixed-points}
Translations have no fixed points: if
$x\in G$ is such that $x=gx$, then
$g=e$.
\end{remark}

\begin{definition}
\label{definition-fixed-point-property}
A topological space $X$ has the
{\it fixed point property} if every
continuous function $f:X\to X$ has a
fixed point.
\end{definition}

Some examples of spaces with the fixed
point property are the interval
$[0,1]\subset\mathbb{R}$, the $n$-balls
$B_r(0)\subset\mathbb{R}^n$ and the
Hilbert cube $[0,1]^{\mathbb{N}}$. Since
translations have no fixed points, the
Hilbert cube cannot have the structure
of a topological group (with its usual
topology, of course).

Now consider a subgroup $H<G$ and the
quotient space $G/H$, which is defined
even if $H$ is not normal (it is a
topological space without group
structure). It turns out that no space
of this form is homeomorphic to the
Hilbert cube.

\noindent
{\it Note on the original.} The balls
should be closed balls. The last claim
is stated without proof or reference.

\begin{remark}
\label{remark-sum-of-closed-sets}
The set
$$
A=\left\{n+\frac{1}{n}
\right\}_{n\in\mathbb{N}}
$$
is closed, but
$$
A-\mathbb{N}
=\left\{\frac{1}{n}
\right\}_{n\in\mathbb{N}}
$$
is not.
\end{remark}

\noindent
{\it Note on the original.} The set
$A-\mathbb{N}$ of all differences is
larger than the set displayed; it
contains the displayed points and has
$0$ as a limit point outside it, which
is the point of the remark: the product
of two closed sets need not be closed.

\begin{theorem}
\label{theorem-products-of-subsets}
Let $G$ be a topological group,
$A,B,O,U\subset G$ and $g\in G$. Then
\begin{enumerate}
\item if $A$ is closed and $B$ is
compact, $A\cdot B$ is closed;
\item if $A$ and $B$ are open,
$A\cdot B$ is open;
\item if $A$ is open (closed),
$g\cdot A$ and $A\cdot g$ are open
(closed);
\item if $A$ and $B$ are compact,
$A\cdot B$ is compact.
\end{enumerate}
\end{theorem}

\begin{proposition}
\label{proposition-coset-projection}
Let $G$ be a topological group and
$H<G$. The projection
$$
p:G\to G/H
$$
is continuous, open and, if $G$ is
compact, closed.
\end{proposition}

\noindent
{\it Note on the original.} This
statement appears there under the
heading ``Definition'' and without
proof. Openness is
Corollary
\ref{corollary-coset-projection-open}.

\section{Group actions}
\label{section-group-actions}

\begin{definition}
\label{definition-action}
Let $G$ be a topological group and $X$
a topological space. An {\it action}
$G\curvearrowright X$ is a continuous
function $\alpha:G\times X\to X$ such
that
\begin{enumerate}
\item for $g,h\in G$ and $x\in X$,
$$
\alpha(g,\alpha(h,x))
=\alpha(g\cdot h,x),
$$
that is, $g(hx)=(g\cdot h)x$;
\item $\alpha(e,x)=x$ where $e\in G$ is
the identity, that is, $ex=x$.
\end{enumerate}
The triple $(G,X,\alpha)$ is called a
{\it $G$-space} or a
{\it topological transformation group}.
\end{definition}

This last name refers to the
transformations of the form
$$
\tilde{g}:X\to X,
\qquad
x\mapsto gx
$$
for $g\in G$, which are continuous.
Moreover,
\begin{enumerate}
\item $\widetilde{gh}
=\tilde{g}\circ\tilde{h}$,
\item $\tilde{e}=1_X$.
\end{enumerate}

\begin{proposition}
\label{proposition-translation-homeomorphism}
The map $\tilde{g}$ is a homeomorphism.
\end{proposition}

\begin{proof}
We have
$$
\tilde{g}\circ\widetilde{g^{-1}}
=\widetilde{gg^{-1}}
=\tilde{e}=1_X,
$$
and also
$\widetilde{g^{-1}}\circ\tilde{g}=1_X$,
so
$\widetilde{g^{-1}}=(\tilde{g})^{-1}$.
\end{proof}

So in fact we have the homomorphism
$$
G\to\mathcal{H}(X),
\qquad
g\mapsto\tilde{g}
$$
where $\mathcal{H}(X)$ is the group of
homeomorphisms of $X$.

The interesting theory appears when $X$
has properties, such as being a linear
space. In that case we could take the
linear homeomorphisms and build a
diagram where we {\it represent} the
group $G$ by linear homeomorphisms.

Now, is the map $g\mapsto\tilde{g}$
continuous? To answer this question we
need to give a topology to
$\mathcal{H}(X)$. Well,
$\mathcal{H}(X)
\subset\mathcal{C}_C(X,X)$, where the
index $C$ stands for the compact-open
topology on the space of continuous
functions of $X$.

\begin{definition}
\label{definition-equivariant}
If $X$ and $Y$ are two $G$-spaces, a
function $f:X\to Y$ is called
{\it $G$-equivariant} if
$f(gx)=gf(x)$.
\end{definition}

\begin{theorem}[For presentation]
\label{theorem-homeo-group-locally-compact}
Let $X$ be a locally compact Hausdorff
and locally connected space. Then
$\mathcal{H}_C(X)$ is a topological
group.
\end{theorem}

\begin{remark}
\label{remark-homeo-group-locally-connected}
On the proof: if we do not ask for
local connectedness, the continuity of
$f\mapsto f^{-1}$ fails. For example,
the space
$$
X=\left\{\frac{1}{n}
\right\}_{n\in\mathbb{N}}
$$
is locally compact but not locally
connected.
\end{remark}

\noindent
{\it Note on the original.} The space
displayed is discrete, hence locally
connected; the counterexamples I know
need a different space. The remark is
kept as written.

\begin{theorem}[For presentation]
\label{theorem-homeo-group-compact}
Let $X$ be compact and Hausdorff. Then
$\mathcal{H}_C(X)$ is a topological
group.
\end{theorem}

Another interesting topic is the
following.

\begin{theorem}[For presentation]
\label{theorem-evaluation-action}
The action
$$
\mathcal{H}(X)\times X\to X,
\qquad
(f,x)\mapsto f(x)
$$
is continuous when $\mathcal{H}(X)$ is
a topological group. That is, $X$ is an
$\mathcal{H}(X)$-space.
\end{theorem}

\begin{remark}
\label{remark-evaluation-not-continuous}
If we take
$$
\mathcal{C}(X,X)\times X\to X,
\qquad
(f,x)\mapsto f(x)
$$
this is not necessarily continuous, for
example if $X=\mathbb{Q}$, since
$\mathbb{Q}$ is not locally compact.
But if $X$ is locally compact, it is.
\end{remark}

\begin{example}
\label{example-actions}
\begin{enumerate}
\item Let $G$ be any topological group
and take $X=G$. We can define the
action in different ways:
\begin{enumerate}
\item $g\cdot x=gx$, the same operation
as that of the group;
\item $g\cdot x=xg^{-1}$, which is not
the same as the previous one if the
group is not abelian;
\item $g\cdot x=gxg^{-1}$.
\end{enumerate}
\item For a subgroup $H<G$ and the set
of cosets
$$
G/H=\{gH:g\in G\},
$$
we have $G\curvearrowright G/H$ given
by $g(xH)=(gx)H$:
$$
\xymatrix{
G\times G/H \ar[r] & G/H \\
G\times G \ar@<0.5ex>[u]
  \ar@<-0.5ex>[u] \ar[r]
  & G \ar[u]_{\pi}
}
$$
\item If $X$ is a $G$-space and $H<G$,
there is an induced action $\alpha'$ of
$H$, the restriction of $\alpha$ to
$H\times X$:
$$
\xymatrix{
G\times X \ar[r]^-{\alpha} & X \\
H\times X \ar@{^{(}->}[u]
  \ar@{-->}[ur]_{\alpha'}
}
$$
\item $G=\mathbb{S}^1$ and
$X=\mathbb{R}^2$, where $G$ acts by
multiplication; geometrically these are
rotations.
\item
$\operatorname{SO}(n)
\curvearrowright\mathbb{R}^n$, which we
now see as orthogonal transformations
that preserve the norm. To see that the
action is continuous, we use the
previous theorems to see that the
evaluation
$$
\omega:
\mathcal{H}_C(\mathbb{R}^n)
\times\mathbb{R}^n
\to\mathbb{R}^n
$$
is continuous, and then use
$$
\xymatrix{
\mathcal{H}_C(\mathbb{R}^n)
  \times\mathbb{R}^n
  \ar[r]^-{\omega}
  & \mathbb{R}^n \\
\operatorname{SO}(n)
  \times\mathbb{R}^n
  \ar@{^{(}->}[u] \ar@{-->}[ur]
}
$$
with the detail that one has to check
that the topology of
$\operatorname{SO}(n)$ is the same as
the compact-open topology.
\item In the previous diagram we can
replace $\operatorname{SO}(n)$ by
$\operatorname{O}(n)$ or
$\operatorname{GL}(n)$.
\end{enumerate}
\end{example}

\noindent
{\it Note on the original.} In the last
diagram the original omits the factors
$\mathbb{R}^n$ on the left.

\begin{definition}
\label{definition-orbit}
Let $X$ be a $G$-space and $x\in X$.
The {\it orbit} of $x$ is
$$
G(x)=\{gx:g\in G\}.
$$
\end{definition}

Note that two orbits are disjoint or
equal. In fact, setting $x\sim y$ when
$G(x)=G(y)$ gives an equivalence
relation, from which we get the
{\it orbit projection}
$\pi:X\to X/G$, and we can give this
set the quotient topology.

\begin{proposition}
\label{proposition-orbit-projection-open}
The orbit projection is open.
\end{proposition}

\begin{proof}
Take $U\subset X$ open. To see that
$\pi(U)$ is open, it is enough to see
that $\pi^{-1}(\pi(U))$ is open. A
point $x\in X$ lies in
$\pi^{-1}(\pi(U))$ if and only if
$\pi(x)\in\pi(U)$, that is, if and only
if there is $u\in U$ with
$\pi(x)=\pi(u)$, that is, with
$x\in G(u)$. Hence
$$
\begin{aligned}
\pi^{-1}\left(\pi(U)\right)
&=\bigcup_{u\in U}G(u) \\
&=\bigcup_{g\in G}gU
\end{aligned}
$$
which is open.
\end{proof}

\begin{corollary}
\label{corollary-coset-projection-open}
The projection $p:G\to G/H$ is open.
\end{corollary}

\begin{definition}
\label{definition-stabilizer}
Let $X$ be a $G$-space. The
{\it stabilizer} of $x\in X$ is the
subgroup
$$
G_x=\{g\in G:gx=x\}<G.
$$
\end{definition}

\begin{proposition}
\label{proposition-stabilizer-closed}
If $X$ is Hausdorff, $G_x$ is closed.
\end{proposition}

\begin{definition}
\label{definition-free-action}
An action is {\it free} if all the
stabilizers are trivial, that is,
$G_x=\{e\}$ for every $x\in X$.
\end{definition}

\begin{definition}
\label{definition-transitive-action}
An action is {\it transitive} if
$G(x)=X$ for every $x\in X$. In this
case $X/G$ is a single point.
\end{definition}

\begin{example}
\label{example-pathological-orbit-spaces}
\begin{enumerate}
\item $X=\mathbb{R}^2$ and
$G=(\mathbb{R}_+,\cdot)$: the action by
rays. This space is $T_0$ but not
$T_1$.
\item We say that
$(\mathbb{R},X,\alpha)$ is a
{\it dynamical system} and that
$(\mathbb{Z},X,\alpha)$ is a
{\it discrete dynamical system}.
\item
$\mathbb{R}\curvearrowright\mathbb{S}^1$
by
$$
t\cdot e^{iu}=e^{i(t+u)}.
$$
\item $\mathbb{R}\curvearrowright
\mathbb{S}^1\times\mathbb{S}^1$ by
$$
t\cdot(e^{iu},e^{iv})
=(e^{i(u+t)},e^{i(u+\sqrt{2}t)}).
$$
The orbits $\mathbb{R}(z)$ are dense in
$\mathbb{S}^1\times\mathbb{S}^1$. This
space is neither $T_1$ nor $T_0$.
\end{enumerate}
\end{example}

\noindent
{\it Note on the original.} In the
first and last items ``this space'' is
the orbit space $X/G$. In the last
formula the second coordinate should be
$e^{i(v+\sqrt{2}t)}$.

\begin{remark}
\label{remark-pullback-action}
For an epimorphism $p:G'\to G$, if $X$
is a $G$-space, it is also a
$G'$-space.
\end{remark}

The following theorem rules out these
pathological cases.

\begin{theorem}
\label{theorem-compact-group-orbits}
Let $X$ be a $G$-space with $G$ a
compact group.
\begin{enumerate}
\item Every orbit is closed in $X$.
\item The projection $p$ is closed.
\item There is a $G$-equivariant
homeomorphism
$$
G(x)\approx_GG/G_x,
$$
since both are $G$-spaces.
\end{enumerate}
\end{theorem}

\begin{proof}
(1) Fix $x\in X$ and define
$$
p:G\to G(x),
\qquad
g\mapsto gx.
$$
Since $p$ is continuous and $G$ is
compact, $G(x)$ is compact and in
particular closed.

(3) We take the action
$G\curvearrowright G(x)$ given by
$g'(gx)=(g'g)x$. This makes the orbit
projection $\eta$ of the action of
$G_x$ on $G$ open. Now observe
$$
\xymatrix{
G \ar[r]^-{p} \ar[d]_{\eta}
  & G(x) \\
G/G_x \ar@{-->}[ur]_{\tilde{p}}
}
$$
Note that $p$ is constant on the
cosets, so we can define the dashed
arrow by
$$
\tilde{p}(gG_x):=p(g)=gx.
$$
This function turns out to be bijective
and, by compactness, closed.

To see that $\tilde{p}$ is equivariant,
$$
\tilde{p}(t\cdot gG_x)
=t\,\tilde{p}(gG_x).
$$
\end{proof}

\noindent
{\it Note on the original.} The second
item is numbered 1.5 there and has no
proof. In the proof of the third item
the original says ``the action
$G(x)\curvearrowright G$'', writes
$\tilde{p}(gG_x)=gG_x$ and states the
equivariance as
$\tilde{p}(t\cdot g)=t\tilde{p}(g)$;
the text above gives the intended
reading. Closedness of the orbits and of
$\tilde{p}$ uses that $X$ is Hausdorff.

The simplest examples of $G$-spaces are
of the form $G(x)=G/H$, in particular
$G(x)=G/G_x$. And they are glued
together like this: $X$ is the disjoint
union of its orbits $G(x)$.

\begin{definition}
\label{definition-invariant-subset}
A subset $A\subset X$ is
{\it $G$-invariant} if
$G(a)\subseteq A$ for every $a\in A$.
We can take restricted actions to
produce {\it $G$-subspaces}:
$$
\xymatrix{
G\times X \ar[r]^-{\alpha} & X \\
G\times A \ar@{^{(}->}[u]
  \ar[r]^-{\alpha'}
  & A \ar@{^{(}->}[u]
}
$$
\end{definition}

\begin{definition}
\label{definition-perfect-map}
A surjective function $f:X\to Y$ is
{\it perfect} if the fibers
$f^{-1}(y)$ are compact.
\end{definition}

\noindent
{\it Note on the original.} A perfect
map is also required to be continuous
and closed, as the original itself says
later, in the proof of
Proposition
\ref{proposition-invariant-refinement}.

\begin{theorem}
\label{theorem-perfect-invariants}
Let $G$ be compact and $X$ a $G$-space.
Then any property $\mathcal{P}$ that is
invariant under perfect maps passes
from $X$ to $X/G$ through the orbit
projection. Among them:
\begin{enumerate}
\item compactness;
\item connectedness;
\item the separation axioms $T_0$,
$T_1$, $T_2$, $T_3$,
$T_{3\frac{1}{2}}$, $T_4$;
\item paracompactness;
\item metrizability.
\end{enumerate}
\end{theorem}

\begin{definition}
\label{definition-invariant-metric}
A metric $\rho$ on a $G$-space $X$ is
{\it invariant} if $\tilde{g}$ is an
isometry for every $g$ in $G$, that is,
$$
\rho(gx,gy)=\rho(x,y).
$$
\end{definition}

In fact, there are times when the
$\tilde{g}$ are linear.

\begin{theorem}
\label{theorem-invariant-metric-quotient}
Let $G$ be compact and $X$ a metrizable
$G$-space. Then there is a compatible
invariant metric on $X/G$.
\end{theorem}

(This theorem implies property 5 of
Theorem
\ref{theorem-perfect-invariants}.)

\begin{proof}
We say that two orbits in a metric
space are {\it parallel} if the
quantity
$$
\rho(x,G(y))
=\inf_{g\in G}\rho(x,gy)
$$
is constant along the orbit of $x$.
This allows us to define a metric on
the orbit space by
$$
\tilde{\rho}(\pi(x),\pi(y))
:=\rho(x,G(y)).
$$
\end{proof}

\noindent
{\it Note on the original.} This is only
the idea; the full argument is the
proof of the next theorem, which assumes
an invariant metric on $X$. That a
metrizable $G$-space with $G$ compact
admits a compatible invariant metric is
not proved in the notes. The statement
presumably means a compatible metric on
$X/G$.

\begin{theorem}[For presentation]
\label{theorem-metric-on-orbit-space}
If $G$ is compact and $X$ is a
$G$-space that admits an invariant
metric, then there is a compatible
metric on $X/G$.
\end{theorem}

\begin{proof}
For fixed $x,y\in X$, the function
$$
G(y)\to\mathbb{R},
\qquad
gy\mapsto\rho(x,gy)
$$
is bounded below by zero, so we can
define
$$
\rho(x,G(y))
:=\inf_{g\in G}\rho(x,gy).
$$
Now we see that in fact this map is
constant on $G(x)$. Since the metric of
$X$ is $G$-invariant, for any
$hx\in G(x)$,
$$
\begin{aligned}
\rho(hx,G(y))
&=\inf_{g\in G}\rho(hx,gy) \\
&=\inf_{g\in G}\rho(x,h^{-1}gy) \\
&=\inf_{g\in G}\rho(x,gy) \\
&=\rho(x,G(y))
\end{aligned}
$$
because the sets
$\{\rho(hx,gy):g\in G\}$ and
$\{\rho(x,h^{-1}gy):g\in G\}$ are
equal, so their infimum is the same.
Define
$$
\tilde{\rho}(G(x),G(y))
:=\rho(x,G(y)).
$$
Now we check that this function is
indeed a metric.
\begin{enumerate}
\item $\tilde{\rho}(G(x),G(x))=0$
trivially.

\item It is positive, that is,
$\tilde{\rho}(G(x),G(y))\geq0$ for any
$x,y\in X$.

If $G(x)\neq G(y)$ but
$\tilde{\rho}(G(x),G(y))=0$, then we
can (since $X$ is metrizable) build a
sequence of points in $G(y)$ whose
distances to some $x\in G(x)$ tend to
zero, that is, the sequence tends to
$x$. Since the group is compact, $G(y)$
is closed and contains its limit
points, so $x\in G(y)$, which is not
possible.

\item It is symmetric, since
$$
\begin{aligned}
\tilde{\rho}(G(x),G(y))
&=\inf_{g\in G}\rho(x,gy) \\
&=\inf_{g\in G}\rho(gy,x) \\
&=\inf_{g\in G}\rho(y,g^{-1}x) \\
&=\inf_{g\in G}\rho(y,gx) \\
&=\tilde{\rho}(G(y),G(x)).
\end{aligned}
$$

\item The triangle inequality is
inherited. Write
$$
m=\inf_{g,h\in G}
\left(\rho(x,hy)+\rho(hy,gz)\right).
$$
Then
$$
\begin{aligned}
\tilde{\rho}(G(x),G(z))
&=\inf_{g\in G}\rho(x,gz) \\
&\leq m \\
&=\inf_{h\in G}\rho(x,hy)
 +\inf_{g,h\in G}\rho(hy,gz) \\
&=\tilde{\rho}(G(x),G(y))
 +\tilde{\rho}(G(y),G(z))
\end{aligned}
$$
since
$$
\rho(x,gz)
\leq\rho(x,hy)+\rho(hy,gz)
$$
for any $g,h\in G$.
\end{enumerate}

Finally, we see that the topology
induced by $\tilde{\rho}$ is the same
as the quotient topology. Recall that
the open sets of the quotient topology
are the subsets of $X/G$ whose preimage
under the projection $\pi:X\to X/G$ is
open in $X$.

{\it Claim.} For any $x\in X$,
$$
\pi(B_\varepsilon(x))
=\tilde{B}_\varepsilon(G(x))
$$
for the open balls of radius
$\varepsilon>0$ in the respective
spaces.

First we see why the result follows
from this claim. Let $\tau$ be the
quotient topology on $X/G$, and take an
open set $U$ of $(X/G,\tau)$. Its
inverse image is an open set in $X$,
which is a union of open balls. The
image of each open ball is an open ball
in $(X/G,\tilde{\rho})$, so
$\pi(\pi^{-1}(U))=U$ is open in
$(X/G,\tilde{\rho})$.

Conversely, any open set of
$(X/G,\tilde{\rho})$ is a union of
balls, each of which is the image of an
open ball in $X$. Since the quotient
projection is open in $(X/G,\tau)$, we
are done.

Now we prove the claim. Take
$\varepsilon>0$, $x\in X$ and the open
balls: $B_\varepsilon(x)$ is the set of
$y\in X$ with $\rho(x,y)<\varepsilon$,
and $\tilde{B}_\varepsilon(G(x))$ is
the set of $G(y)\in X/G$ with
$\tilde{\rho}(G(x),G(y))<\varepsilon$.

($\subseteq$) If
$y\in B_\varepsilon(x)$, then
$\pi(y)=G(y)
\in\tilde{B}_\varepsilon(G(x))$ since
$$
\tilde{\rho}(G(x),G(y))
\leq\rho(x,y).
$$

($\supseteq$) If
$G(y)\in\tilde{B}_\varepsilon(G(x))$,
we must find
$y'\in B_\varepsilon(x)$ such that
$\pi(y')=G(y)$. The distance between
$x$ and the set $\pi^{-1}(G(y))$ is
$$
\begin{aligned}
\rho(x,\pi^{-1}(G(y)))
&=\inf_{z\in\pi^{-1}G(y)}\rho(x,z) \\
&=\inf_{g\in G}\rho(x,gy) \\
&=\rho(x,G(y)) \\
&=\tilde{\rho}(G(x),G(y)) \\
&<\varepsilon.
\end{aligned}
$$
So there is
$y'\in\pi^{-1}(G(y))
\cap B_\varepsilon(x)$.
\end{proof}

The theorem is also true if $G$ is
countably compact.

\begin{definition}
\label{definition-countably-compact}
A space $X$ is {\it countably compact}
if every countable open cover has a
finite subcover.
\end{definition}

To prove this variant of the theorem it
is enough to check that the metric is
positive (the second item of the
proof). It will turn out that the orbit
$G(y)$, being the image of $G$ under
$g\mapsto gy$, is countably compact, so
it must contain at least one
{\it $\omega$-limit point}, which gives
a contradiction as in the previous
proof. Specifically:

\begin{proposition}
\label{proposition-image-countably-compact}
The image of a countably compact space
under a continuous function is
countably compact.
\end{proposition}

\begin{proof}
Suppose that $f:X\to Y$ is continuous
and surjective and that $X$ is
countably compact. Take a countable
cover $\{U_\alpha\}$ of $Y$. Then
$\{f^{-1}(U_\alpha)\}$ is a countable
cover of $X$, for which there is a
finite subcover
$f^{-1}(U_1),\ldots,f^{-1}(U_n)$. If
$U_1,\ldots,U_n$ does not cover $Y$,
choose a point $y$ of $Y$ outside
$U_1\cup\ldots\cup U_n$, and then
$y'\in f^{-1}(y)$, so that
$y'\in f^{-1}(U_j)$ for some
$j=1,\ldots,n$. But then
$f(y')=y\in U_j$.
\end{proof}

\begin{definition}
\label{definition-omega-limit-point}
Let $X$ be a topological space and
$A\subset X$. We say that $x\in X$ is
an {\it $\omega$-limit point} of $A$ if
every open set $U\subset X$ that
contains $x$ has infinitely many points
of $A$.
\end{definition}

In our case, the space $X$ is the orbit
$G(y)$, and the subset $A$ is the
sequence converging to $x$. We will see
that in fact this point is in $A$.

\begin{proposition}
\label{proposition-omega-limit-point}
Let $X$ be a topological space and
$A\subset X$ an infinite subset. If $X$
is countably compact, then $A$ has at
least one $\omega$-limit point.
\end{proposition}

\begin{proof}
(Taken from
\href{https://math.stackexchange.com/%
questions/1817245/%
a-topological-space-is-countably-%
compact-iff-every-countably-infinite-%
subset-has}{this Mathematics Stack
Exchange question}.) Suppose that $A$
has no $\omega$-limit points. Then for
any $x\in X$ there is a neighborhood
$U_x$ whose intersection with $A$ is
finite. Clearly the family of all $U_x$
with $x\in X$ is a cover of $X$. We
will use that the family of finite
subsets of a countable set is countable
to build a countable cover. For any
finite $F\subset A$, let $U_F$ be the
union of the sets $U_x$ such that
$$
U_x\cap A=F.
$$
The sets $U_F$ form a countable cover
of $X$, since they contain all the
points of the $U_x$. Note that every
$U_F\cap A$ is finite by definition.
But if there were a finite subcover,
there would have to be at least one
$U_F$ containing infinitely many points
of $A$.
\end{proof}

\noindent
{\it Note on the original.} The
countability of the family of the $U_F$
needs $A$ to be countable; one first
replaces $A$ by a countably infinite
subset. The sets $U_x$ should be open.

\begin{example}[For presentation]
\label{example-noncompact-metric-fails}
If $G$ is not compact, we can give a
counterexample: a space $X$ with a
$G$-invariant metric that does not pass
to the quotient $X/G$.

It is necessary that the orbits are not
compact sets. The example is the
irrational flow on the torus: since the
orbits are dense, in any neighborhood
of a point $x$ there is an element of
any orbit $G(y)$, so
$$
\inf_{g\in G}\rho(x,gy)=0
$$
even though $G(x)\neq G(y)$.

A simpler action with this idea is
$\mathbb{Z}\curvearrowright
\mathbb{S}^1$ by
$$
n\cdot e^{2\pi it}=e^{2\pi itn}.
$$
It will turn out that the orbit of any
point with irrational argument is dense
in $\mathbb{S}^1$. To prove it we fix
$\alpha$ irrational and obtain from
Dirichlet's approximation theorem that
$$
\begin{aligned}
|e^{2\pi i\alpha n}-e^{2\pi im}|
&\leq|2\pi i\alpha n-2\pi im| \\
&=2\pi|\alpha n-m| \\
&<\frac{1}{N}
\end{aligned}
$$
for any integer $N$, so the orbits
$\mathbb{Z}(\alpha)$ and
$\mathbb{Z}(1)$ are infinitely close.
But the neutral element of $\mathbb{Z}$
sends every point to $1$! {\bf This is
not an action.}

Maybe the multiplicative group
$\mathbb{Q}^\times$ acting on
$\mathbb{R}$?
\end{example}

\noindent
{\it Note on the original.} The second
part of this example is a failed
attempt and is recognized as such
there; the last question is left open.

\begin{example}
\label{example-one-sided-invariant-metric}
If the group is metrizable and compact
then it can always be done. But if the
group is not compact and not abelian it
is possible to give a metric that is
invariant on one side but not on the
other. (See Tkachenko et al.)
\end{example}

\noindent
{\it Note on the original.} The
reference ``Tkachenko et al.'' is not
specified further.

\section{Twisted product}
\label{section-twisted-product}

Let $G$ be a topological group, $H<G$
and $X$ an $H$-space. Then we have the
action $H\curvearrowright(G\times X)$
given by
$$
h\cdot(g,x):=(gh^{-1},hx).
$$
We can then consider the orbit space
$$
G\times_HX:=(G\times X)/H,
$$
which is called the
{\it twisted product} of $G$ and $X$
by means of $H$.

In this space the action of $H$ is
lost, but the action of $G$ is
preserved. We denote the orbit
$H((g,x))$ by $[g,x]$.

\begin{theorem}
\label{theorem-action-on-twisted-product}
The correspondence
$G\curvearrowright(G\times_HX)$ defined
by
$$
g'*[g,x]:=[g'\cdot g,x]
$$
is an action.
\end{theorem}

\begin{proof}
Consider the following diagram, where
$\beta(g',g,x)=(g'\cdot g,x)$ and
$\alpha$ is the action of the
statement:
$$
\xymatrix{
G\times(G\times_HX)
  \ar[r]^-{\alpha}
  & G\times_HX \\
G\times G\times X
  \ar[u]^{\operatorname{id}\times p}
  \ar[r]^-{\beta}
  & G\times X \ar[u]_{p}
}
$$
\end{proof}

\noindent
{\it Note on the original.} The proof
consists of the diagram alone.

Note that we can identify $X$ with the
subset $S$ of $G\times_HX$ of the
classes $[e,x]$ with $x\in X$. The set
$H*S$ of the points $h*[e,x]$ with
$h\in H$ is equal to $S$, since
$$
h*[e,x]=[he,x]
=[heh^{-1}e,x]=[e,x]
$$
because it is an equivalence class.

\noindent
{\it Note on the original.} The last
computation is not right: in general
$$
h*[e,x]=[h,x]=[e,hx],
$$
which is again a point of $S$, so the
conclusion $H*S=S$ holds, but $h$ does
not fix $[e,x]$.

\begin{theorem}
\label{theorem-embedding-twisted-product}
Let $X$ be a topological space. The
function
$$
f:X\to G\times_HX,
\qquad
x\mapsto[e,x]
$$
is an $H$-embedding, that is, an
$H$-equivariant embedding. And it is
closed if $H$ is closed.
\end{theorem}

\noindent
{\it Note on the original.} Here $X$ is
an $H$-space, as in the definition of
the twisted product.

For a proof, we first see the
following.

\begin{proposition}
\label{proposition-closed-set-in-product}
Let $B\subset X$ be a closed set in a
$G$-space $X$, with $G$ any topological
group and $H<G$ a closed subgroup.
Then the set
$$
A=\{(h^{-1},hb):h\in H,b\in B\}
$$
is closed in $G\times X$.
\end{proposition}

\begin{proof}
To prove that it is a closed set we
take a {\it net} (which generalizes a
sequence) and see that it converges.
Let $(h_i)\subset G$ be a net. Suppose
that
$$
(h_i^{-1},h_ib)\to(g,x).
$$
If $h_i^{-1}\to g$ and $h_ib\to x$,
then $h_i^{-1}(h_ib)\to gx$. Then
$gx=b$, so $x=g^{-1}b$ and
$$
(g,x)=(g,g^{-1}b),
$$
and in fact $g\in H$ since $H$ is
closed.
\end{proof}

\noindent
{\it Note on the original.} The net
should be in $H$, and the point $b$
should vary with the index: for a net
$(h_i^{-1},h_ib_i)$ in $A$ converging
to $(g,x)$ one gets $b_i\to gx$, and
$gx\in B$ because $B$ is closed.

\begin{proof}[Proof of Theorem
\ref{theorem-embedding-twisted-product}]
To see that $f$ is injective, consider
the composition
$$
\xymatrix{
X \ar[r] & G\times X \ar[r]^-{\pi}
  & G\times_HX
}
$$
which sends $x$ to $(e,x)$ and then to
$[e,x]$, and $y$ to $(e,y)$ and then to
$[e,y]$. Suppose that $[e,x]=[e,y]$.
Then there is $h\in H$ such that
$$
(e,x)=h\cdot(e,y)=(h^{-1},hy),
$$
but then $h=e$.

To see that $f$ is closed, take
$B\subset X$ a closed set. Then
$p^{-1}(f(B))$ is the set of
$(g,x)\in G\times X$ such that
$[g,x]\in f(B)$, that is, such that
$[g,x]=[e,b]$ for some $b\in B$. So
$$
p^{-1}\left(f(B)\right)
=\{(h^{-1},hb):b\in B,h\in H\}.
$$
Moreover, $f$ is $H$-equivariant, since
$$
f(hx)=[e,hx]=[h,x]=h*[e,x]=hf(x).
$$
\end{proof}

\noindent
{\it Note on the original.} The
equivariance computation is left
unfinished there: it reads
``$[e,hx]=f(hx)=hf(x)=[\,,]$''. The
middle steps above complete it.

\begin{proof}[Proof by Lucho that $f$
is closed]
{\it Claim.}
$\pi^{-1}(S)=H\times X$.

($\supseteq$) Clear, since for
$(h,x)\in H\times X$ we have
$[h,x]=[e,x]$.

($\subseteq$) Let
$(g,x)\in\pi^{-1}(S)$. Then there is
$y\in X$ such that $[g,x]=[e,y]$, so
$$
(g,x)=h\cdot(e,y)=(h^{-1},hy)
$$
for some $h\in H$.

From here it follows that $S$ is
closed.

To see that $f$ is open we build its
inverse through the transgression
theorem applied to
$$
\xymatrix{
H\times X \ar[r]^-{\pi}
  \ar[d]_{\alpha}
  & S \ar@{-->}[dl] \\
X
}
$$
where $\alpha(h,x)=hx$. It is enough to
note that
$$
\pi^{-1}([e,x])
=\{(h,h^{-1}x):h\in H\},
$$
so the dashed arrow sends
$[e,x]\mapsto x$.
\end{proof}

\noindent
{\it Note on the original.} In the
inclusion ($\supseteq$) the class
$[h,x]$ equals $[e,hx]$, not $[e,x]$;
it is still a point of $S$.

\begin{proposition}
\label{proposition-twisted-orbit-space}
For $H<G$ closed there is a
homeomorphism
$$
\frac{G\times_HX}{G}\approx X/H.
$$
\end{proposition}

\begin{proof}
We claim that the map $f$ sending
$H(x)$ to $G([e,x])$ is a
homeomorphism. Since the projection
$p:G\times X\to X$ is $H$-equivariant,
we have
$$
\tilde{p}:G\times_HX\to X/H,
\qquad
[g,x]\mapsto H(x).
$$
It is possible to factor $\tilde{p}$
through $\varphi$ according to
$$
\xymatrix{
\frac{G\times_HX}{G}
  \ar[rd]^{\varphi} \\
G\times_HX \ar[u]^{\pi}
  \ar[r]_-{\tilde{p}}
  & X/H
}
$$
and in fact $\varphi$ and $f$ are
inverse to each other.
\end{proof}

\noindent
{\it Note on the original.} There the
map $f$ is written
$H(x)\mapsto G(x)$ and the target of
$\tilde{p}$ is written $X$.

\section{Slices}
\label{section-slices}

\begin{definition}
\label{definition-cross-section}
A {\it cross section} is a function
$\sigma$ such that the following
diagram commutes:
$$
\xymatrix{
& X \ar[d]^{\pi} \\
X/G \ar[r]_-{\operatorname{Id}}
  \ar@{-->}[ur]^{\sigma}
  & X/G
}
$$
\end{definition}

\begin{definition}
\label{definition-retraction}
Given a subset $A\subseteq X$ of a
topological space, a {\it retraction}
is a function $r:X\to A$ such that
$r\circ i=\operatorname{Id}_A$ for the
inclusion $i:A\hookrightarrow X$.
\end{definition}

\begin{definition}
\label{definition-slice}
Let $X$ be a $G$-space, $H<G$ and
$S\subset X$. We say that $S$ is an
{\it $H$-slice} of $X$ if
\begin{enumerate}
\item $S$ is $H$-invariant, that is,
$H(S)=S$;
\item $S$ is closed in its
{\it saturation}
$$
G(S)=\{gs:g\in G,s\in S\};
$$
\item if $g$ is such that
$gS\cap S\neq\varnothing$, then
$g\in H$;
\item $G(S)$ is open in $X$.
\end{enumerate}
If $G(S)=X$ we say that $S$ is a
{\it global slice}.
\end{definition}

It follows that if two orbits have
nonempty intersection, then they are
the same.

When do slices exist?

\begin{theorem}[Slice theorem]
\label{theorem-slice-first}
If $G$ is a compact Lie group and $X$
is Tychonoff, then for any $x\in X$
there is a $G_x$-slice $S\subset X$
with $x\in S$.
\end{theorem}

\begin{proposition}
\label{proposition-action-on-slice-open}
Let $X$ be a $G$-space with $G$
compact. If $S\subset X$ is a global
$H$-slice and $\alpha$ is the action of
$G$ on $X$, then the restriction of
$\alpha$
$$
G\times S\to G(S)=X
$$
is open.
\end{proposition}

\begin{remark}
\label{remark-slice-of-saturation}
$S$ is a global $H$-slice of $G(S)$, so
for any open sets $U$ of $S$ and $O$ of
$G$, the set $U\cdot O$ is open in
$G(S)$. Moreover, since $G(S)$ is open
in $X$, so is $U\cdot O$.
\end{remark}

\begin{proposition}
\label{proposition-small-slices}
Suppose that $(X,d)$ is a metrizable
$G$-space with $G$ a compact Lie group.
Then for each $x\in X$ and $r>0$ there
is a $G_x$-slice $S$ such that
$$
\operatorname{diam}(S)\leq2r.
$$
\end{proposition}

\begin{definition}
\label{definition-linear-g-space}
We say that $V$ is a
{\it linear $G$-space} if the action is
linear, that is,
$$
g\cdot(\lambda x+y)
=\lambda gx+gy.
$$
\end{definition}

\begin{example}
\label{example-slices}
\begin{enumerate}
\item $S=\{eH\}$ with $eH\in G/H$ is an
$H$-slice. The motion of $eH$ in $G/H$
``models'' the motion of $S$ in $X$.
\item In the twisted product
$G\times_HS$, the subset of the classes
$[e,s]$ with $s\in S$, which we
identify with $S$, is a global
$H$-slice. Let us look at condition 2.
Suppose that
$gS\cap S\neq\varnothing$, so that
there is $h\in H$ such that
$$
[e,x]=[g^{-1},gs]=[h^{-1},hs].
$$
\end{enumerate}
\end{example}

\noindent
{\it Note on the original.} The
condition being checked is the third
one of
Definition \ref{definition-slice}, and
the verification stops there.

\begin{proposition}
\label{proposition-retract-of-orbit-space}
If $H<G$ is closed, the space
$\frac{G\times_HX}{G}$ is homeomorphic
to a retract of $\frac{G\times_HX}{H}$.
\end{proposition}

\begin{proof}
We use that
$\frac{G\times_HX}{G}\approx X/H$
(Proposition
\ref{proposition-twisted-orbit-space}).
Define the retraction
$$
r:\frac{G\times_HX}{H}\to X/H,
\qquad
H([g,x])\mapsto H(x).
$$
Now for the inclusion
$i:X\hookrightarrow G\times_HX$ we have
$$
\tilde{i}:X/H
\hookrightarrow\frac{G\times_HX}{H},
\qquad
H(x)\mapsto H([e,x]),
$$
and $r$ sends $H([e,x])$ back to
$H(x)$, so we have built a left inverse
of $r$, hence it is a retraction
according to the following proposition.
\end{proof}

\noindent
{\it Note on the original.} The map
$\tilde{i}$ is a right inverse of $r$:
the composition $r\circ\tilde{i}$ is
the identity.

\begin{proposition}
\label{proposition-retraction-onto-image}
Let $r:X\to A$ and $f:A\to X$ be
continuous and such that
$r\circ f=\operatorname{Id}_A$. Then
$r:X\to f(A)$ is a retraction onto
$f(A)$.
\end{proposition}

\begin{proof}
The map $rf$ is injective, since if
$rf(a_1)=rf(a_2)$, then $a_1=a_2$.
Moreover, $f$ is an embedding, since
$$
f(C)=f(A)\cap r^{-1}(C)
$$
and $C=r(f(C))$.
\end{proof}

\noindent
{\it Note on the original.} The
retraction onto $f(A)$ is the
composition $f\circ r$.

Now we see that any tubular
neighborhood (the saturation of a
slice) is a twisted product.

\begin{theorem}
\label{theorem-tube-is-twisted-product}
Let $G$ be a topological group, $H<G$
closed, $X$ a $G$-space and
$S\subset X$ an $H$-slice. Then
$$
G\times_HS\approx_GG(S).
$$
\end{theorem}

\begin{proof}
Take
$$
f:G\times_HS\to G(S),
\qquad
[g,s]\mapsto gs.
$$
First note that $f$ sends
$[gh^{-1},hs]$ to
$$
(gh^{-1})hs=gs,
$$
so $f$ is well defined. Then
$$
\xymatrix{
G\times_HS \ar[r] & G(S) \\
G\times S \ar[u] \ar[ur]_{\varphi}
  \ar@{^{(}->}[r]
  & G\times G(S) \ar[u]
}
$$
where $\varphi$ is the action. By this
diagram, $f$ is continuous. Now we see
that it is injective, taking $[g,s]$
and $[g',s']$ that go to $gs=g's'$.
Then
$$
s'=(g')^{-1}gs
\in S\cap(g')^{-1}gS
\neq\varnothing.
$$
Take $h:=(g')^{-1}g\in H$, so that
$g'=gh^{-1}$. It is also surjective.

(I think that) if $G$ is compact, then
$f$ is closed.
\end{proof}

\noindent
{\it Note on the original.} That the
inverse of $f$ is continuous is not
proved; the last sentence is a guess
recorded as such in the original.

\begin{lemma}
\label{lemma-compact-action-closed}
Let $G$ be a compact group and $X$ a
$G$-space. Then the action
$G\times X\to X$ is closed.
\end{lemma}

\begin{theorem}[Slice theorem]
\label{theorem-slice}
Let $G$ be a compact Lie group, $X$ a
Tychonoff $G$-space and $x\in X$. Then
there is a $G_x$-slice $S\subset X$
such that $x\in S$.
\end{theorem}

\begin{proposition}
\label{proposition-preimage-of-slice}
Let $G$ be any topological group and
$f:X\to Y$ a $G$-map. If $S\subset Y$
is an $H$-slice in $Y$, then
$f^{-1}(S)$ is an $H$-slice in $X$.
\end{proposition}

\section{Extension of equivariant maps:
the Tietze--Urysohn theorem}
\label{section-extension}

\noindent
{\it Note on the original.} The title
there reads ``teorema de
Tits-{\it Urison}''; I take it to be the
Tietze--Urysohn extension theorem,
stated below as
Theorem \ref{theorem-tietze}.

Given the action
$\mathbb{Z}_2\curvearrowright\mathbb{R}$
given by $x\mapsto-x$, how do we extend
equivariantly a function from a closed
set $A$ in $\mathbb{R}$ to
$\mathbb{R}^2$?
$$
\xymatrix{
A \ar[r]^-{\varphi}
  \ar@{^{(}->}[d]
  & \mathbb{R} \\
\mathbb{R}^2 \ar@{-->}[ur]_{f}
}
$$
We simply take the average
$$
f(x)=\frac{\varphi(x)-\varphi(-x)}{2}
$$
since in this case $f(x)=-f(x)$, that
is, $f$ is $\mathbb{Z}_2$-equivariant.

\noindent
{\it Note on the original.} The
equivariance is $f(-x)=-f(x)$. The
formula uses $\varphi$ on all of the
space, so $\varphi$ here should already
be some continuous extension of the
given function on $A$ (which exists by
Theorem \ref{theorem-tietze}), with $A$
invariant; the average then makes it
equivariant. The same applies to the
integral below. It is also not clear
whether the domain is $\mathbb{R}$ or
$\mathbb{R}^2$.

Now consider the action
$\mathbb{S}^1
\curvearrowright\mathbb{R}^2$ by
rotations. Take
$$
f(x)
=\int_{\mathbb{S}^1}
g^{-1}\varphi(gx)\,dg
$$
where $dg$ is the {\it Haar measure},
which is defined as follows: it is a
map
$$
\int:\mathcal{C}(G,\mathbb{R})
\to\mathbb{R}
$$
such that
\begin{enumerate}
\item it is linear,
\item $\int1=1$,
\item $\int f\geq0$ if $f\geq0$,
\item $\int gf=\int f=\int g*f$, where
$gf$ and $g*f$ are the left and right
translates of $f$ by $g$.
\end{enumerate}
Now we see whether $f$ is equivariant,
that is, whether $f(hx)=hf(x)$. On the
one hand, with the change of variable
$t=gh$, for which $dt=dg$,
$$
\begin{aligned}
f(hx)
&=\int g^{-1}\varphi(ghx)\,dg \\
&=\int ht^{-1}\varphi(tx)\,dt \\
&=h\int t^{-1}\varphi(tx)\,dt.
\end{aligned}
$$
On the other hand
$$
hf(x)=h\int g^{-1}\varphi(gx)\,dg.
$$
If $a\in A$,
$$
\begin{aligned}
f(a)
&=\int g^{-1}\varphi(ga)\,dg \\
&=\int g^{-1}g\varphi(a)\,dg \\
&=\int\varphi(a)\,dg \\
&=\varphi(a).
\end{aligned}
$$

\noindent
{\it Note on the original.} The domain
of the integral is written
``$(G,\mathbb{R})$'' there, and the
last line of the computation of $f(hx)$
reads
``$h\varphi t^{-1}\varphi(tx)$''. The
last computation uses that $\varphi$ is
equivariant on the invariant set $A$.

\section{Dugundji's theorem}
\label{section-dugundji}

\begin{theorem}[Dugundji]
\label{theorem-dugundji}
Let $X$ be a metrizable space,
$A=\bar{A}\subset X$, let $L$ be a
locally convex vector space and
$V\subset L$ convex. Then any
continuous function $f:A\to V$ admits
an extension $F$.
$$
\xymatrix{
A \ar[r]^-{f} \ar@{^{(}->}[d]
  & V \\
X \ar@{-->}[ur]_{F}
}
$$
Moreover
$$
\operatorname{img}F
\subset\operatorname{conv}
\operatorname{img}F.
$$
\end{theorem}

\noindent
{\it Note on the original.} The last
inclusion is trivial as written; the
statement of Dugundji's theorem is that
the image of $F$ lies in the convex
hull of the image of $f$.

First we will need the following.

\begin{lemma}
\label{lemma-dugundji-system}
In this situation there is an open
cover $\omega$ of $X\setminus A$ such
that
\begin{enumerate}
\item $\omega$ is locally finite in
$X\setminus A$, that is, every point
has a neighborhood that intersects
finitely many open sets of $\omega$;
\item for any point $a\in A$ and each
neighborhood $U_a$ there is another
neighborhood $V_a$ such that if
$V_a\cap W\neq\varnothing$ for some
$W\in\omega$, then $W\subset U_a$.
\end{enumerate}
We say that this cover is a
{\it Dugundji system}.
\end{lemma}

\begin{definition}
\label{definition-partition-of-unity}
A {\it partition of unity} subordinate
to a cover $\omega=\{U_i\}_{i\in J}$ is
a family of continuous functions
$\varphi_i:X\to[0,1]$ with $i\in J$
such that
\begin{enumerate}
\item
$\varphi_i^{-1}((0,1])\subset U_i$,
\item $\sum_{i\in J}\varphi_i(x)=1$.
\end{enumerate}
\end{definition}

It turns out that any locally finite
cover admits a subordinate partition of
unity. Now we prove Dugundji's theorem.

\begin{proof}[Proof of Theorem
\ref{theorem-dugundji}]
Let $\omega$ be the Dugundji system for
$X\setminus A$. For each
$U_i\in\omega$, let $x_i\in U_i$ and
$a_i\in A$ be such that
$$
\rho(x_i,a_i)\leq2d(x_i,A)
$$
where
$$
d(x_i,A):=\inf_{a\in A}\rho(x_i,a).
$$
For a partition of unity of $\omega$,
define $F(x)=f(x)$ if $x\in A$ and
$$
F(x):=\sum_{i\in J}\varphi_i(x)f(a_i)
$$
if $x\in X\setminus A$. Since this is a
partition of unity, the sum is a convex
sum (the coefficients add up to 1), so
in either case $F(x)\in V$.

Now we see that $F$ is continuous. We
proceed by cases.
\begin{enumerate}
\item If $x\in X\setminus A$, there is
$S$ with $x\in S$ such that
$S\cap U_i\neq\varnothing$ for finitely
many indices $i\in J$. Then
$$
F(z)=\sum_{i\in J}\varphi_i(z)f(a_i)
$$
for every $z\in S$ because the
$\varphi_i$ are a partition of unity.

\item If $x=a\in A$, then $F(a)=f(a)$.
Take a convex open set $Q\subset V$
that contains this point. Since $f$ is
continuous on $A$, there is a
neighborhood of $a$ in $A$, that is,
there is $\delta>0$ such that
$$
f(B(a,\delta)\cap A)\subset Q.
$$

By property (2) in
Lemma \ref{lemma-dugundji-system}, we
know that there is $r>0$ such that for
any $W\in\omega$ with
$W\cap B(a,r)\neq\varnothing$ we have
$W\subset B(a,\delta)$. Suppose that
$r<\frac{\delta}{4}$.

We claim that $F(B(a,r))\subset Q$. Let
$x\in B(a,r)\cap(X\setminus A)$. Then
$$
F(x)=\sum_{i\in J}\varphi_i(x)f(a_i)
$$
and we ask whether it lies in $Q$. We
have
$$
\rho(a,a_i)
\leq\rho(a_i,x_i)+\rho(x_i,a)
<2d+\delta<4\delta<r.
$$
Then $a_i\in B(a,r)$, so $f(a_i)\in Q$
and by convexity
$$
F(x)=\sum_{i\in J}\varphi_i(x)f(a_i)
\in Q.
$$
\end{enumerate}
\end{proof}

\noindent
{\it Note on the original.} The chain of
inequalities at the end does not make
sense as written (it ends with
$4\delta<r$ while $r<\delta/4$), and
$d$ is not defined. The standard
argument is as follows: only the
indices $i$ with $x\in U_i$ contribute
to $F(x)$; for those,
$U_i\subset B(a,\delta/3)$ if $r$ is
chosen for the neighborhood
$B(a,\delta/3)$, so
$\rho(x_i,a)<\delta/3$,
$\rho(x_i,a_i)\leq2d(x_i,A)
<2\delta/3$ and
$\rho(a,a_i)<\delta$. In the first case
the set $S$ should be an open
neighborhood of $x$ inside
$X\setminus A$.

Now we ask what happens in the
equivariant analogue.

\begin{definition}
\label{definition-absolute-extensor}
Let $X$ be a topological space and $K$
some class of topological spaces. $X$
is said to be an
{\it absolute extensor} for $K$,
written $X\in\operatorname{AE}(K)$, if
for any $Y\in K$, any closed
$A\subset Y$ and any continuous
$f:A\to X$ there is a continuous
extension $F:Y\to X$.
\end{definition}

\begin{example}
\label{example-convex-absolute-extensor}
For $W$ a convex subset of a locally
convex space and $M$ the class of
metrizable spaces, Dugundji's theorem
ensures that
$W\in\operatorname{AE}(M)$.
\end{example}

\begin{theorem}[Tietze]
\label{theorem-tietze}
Let $X$ be a normal space and
$A\subset X$ closed. If
$f:A\to[0,1]$ is a continuous function,
then there is a continuous extension
$f:X\to[0,1]$.
\end{theorem}

\begin{definition}
\label{definition-neighborhood-extensor}
Let $X$ be a topological space and $K$
some class of topological spaces. $X$
is said to be an
{\it absolute neighborhood extensor}
for $K$, written
$X\in\operatorname{ANE}(K)$, if for any
$Y\in K$, any closed $A\subset Y$ and
any continuous $f:A\to X$ there is a
continuous extension $F:U\to X$ with
$U$ an open neighborhood of $A$.
\end{definition}

\noindent
{\it Note on the original.} In both
definitions the original writes
``$y\in K$, $A\in Y$'' and does not say
that $A$ is closed or that $F$ extends
$f$.

\begin{theorem}
\label{theorem-ae-iff-contractible-ane}
$X\in\operatorname{AE}(M)$ if and only
if $X$ is contractible and
$X\in\operatorname{ANE}(M)$.
\end{theorem}

\begin{theorem}[Dugundji for $G$-spaces]
\label{theorem-equivariant-dugundji}
Let $X$ be a metrizable $G$-space and
$A=\bar{A}\subset X$ invariant, that
is, $G(A)=A$. If $L$ is a locally
convex vector $G$-space and
$V\subset L$ is invariant and convex,
then any equivariant function
$f:A\to V$ admits an extension
$F:U\to V$ with $U$ an invariant
neighborhood of $A$.
$$
\xymatrix{
A \ar[r]^-{f} \ar@{^{(}->}[d]
  & V \\
U \ar@{-->}[ur]_{F}
}
$$
Moreover
$$
\operatorname{img}F
\subset\operatorname{conv}
\operatorname{img}F.
$$
\end{theorem}

\noindent
{\it Note on the original.} The
hypothesis reads there ``any function
$f:A\to V$ is equivariant, then
admits\ldots''. The extension $F$ is
presumably meant to be equivariant, and
the last inclusion has the same problem
as in Theorem \ref{theorem-dugundji}.
No hypothesis on $G$ is stated.

That is, $f$ cannot be extended to the
whole space but it can be extended to a
neighborhood. It will turn out that if
$X$ is contractible we can extend the
function. Giving the spaces of
continuous functions the compact-open
topology, they turn out to be
topological vector spaces too (the
operations are continuous) and
$G$-spaces (with
$(gf)(x)=g^{-1}f(gx)$), and the
operator
$$
\Delta:\mathcal{C}(A,V)
\to\mathcal{C}(X,V)
$$
is a linear embedding.

\noindent
{\it Note on the original.} The action
on functions is written
``$gf(x)=g^{-1}f(gx(x))$'' there; the
usual convention is
$(gf)(x)=gf(g^{-1}x)$.

\section{$G$-canonical cover and
Dugundji's theorem}
\label{section-canonical-cover}

\begin{definition}
\label{definition-refinement-paracompact}
\begin{enumerate}
\item Let $\mathcal{U},\mathcal{V}$ be
open covers of a topological space $X$.
We say that $\mathcal{U}$
{\it refines} $\mathcal{V}$ if for any
$U\in\mathcal{U}$ there is
$V\in\mathcal{V}$ such that
$U\subset V$.
\item Let $\mathcal{U}$ be an open
cover of $X$. We say that $\mathcal{U}$
is {\it locally finite} if any
$x\in X$ has a neighborhood that
intersects finitely many elements of
$\mathcal{U}$.
\item A space is {\it paracompact} if
any open cover admits a locally finite
refinement.
\end{enumerate}
\end{definition}

\begin{remark}
\label{remark-compact-paracompact}
If $X$ is compact then it is
paracompact.
\end{remark}

\begin{theorem}[Stone]
\label{theorem-stone}
If $X$ is metrizable then $X$ is
paracompact.
\end{theorem}

\begin{proof}
Similar to that of Urysohn's theorem.
\end{proof}

So if $X$ is metrizable and the metric
is $G$-invariant, then $X/G$ is
paracompact.

\begin{theorem}
\label{theorem-paracompact-partitions}
$X$ is paracompact if and only if any
open cover $\mathcal{U}$ of $X$ admits
a partition of unity subordinate to
$\mathcal{U}$.
\end{theorem}

\begin{proposition}
\label{proposition-map-from-fixed-points}
Let $G$ be a topological group,
$H\leq G$ closed and $X$ a $G$-space.
Let $X[H]$ be the set
$\operatorname{Fix}(H)$ of the points
$x\in X$ such that $hx=x$ for all
$h\in H$. The function
$$
G/H\times X[H]\to X,
\qquad
(gH,x)\mapsto gx
$$
is well defined and continuous.
\end{proposition}

\begin{proof}
Let $(g_1H,x)$ and $(g_2H,x)$ be points
of $G/H\times X[H]$ such that
$g_1H=g_2H$, that is,
$g_2^{-1}g_1\in H$. Then
$g_2^{-1}g_1x=x$, so $g_1x=g_2x$.

To see that it is continuous, simply
look at
$$
\xymatrix{
G/H\times X[H] \ar[r]^-{\omega}
  & X \\
G\times X[H] \ar[ur]_{\alpha}
  \ar[u]^{p\times\operatorname{Id}}
}
$$
where $\alpha$ is the restriction of
the action to $G\times X[H]$. Assuming
for now that $X[H]$ is closed, the
preimages of closed sets under $\omega$
are closed, since
$p\times\operatorname{Id}$ is open.
\end{proof}

\begin{definition}
\label{definition-invariant-cover}
Let $X$ be a $G$-space. We say that a
cover $\{U_\alpha\}$ is
{\it invariant} if each $U_\alpha$ is
an invariant set (that is, a tube: the
saturation of a slice).
\end{definition}

\begin{proposition}
\label{proposition-invariant-refinement}
Let $X$ be a paracompact space. If
$\{U_\alpha\}$ is an invariant cover,
then there is a locally finite
refinement of $\{U_\alpha\}$ by
invariant sets.
\end{proposition}

\begin{proof}
Since the projection is a perfect map
(continuous, closed and with compact
fibers), $X/G$ is paracompact.

Let $\{U_\alpha\}$ be an invariant
cover of $X$. Then $\{p(U_\alpha)\}$ is
an open cover of $X/G$, so we can find
a locally finite refinement
$\{V_\lambda\}$. Then
$\{p^{-1}(V_\lambda)\}$ is an open
cover of $X$; in fact, it is locally
finite since $\{V_\lambda\}$ is locally
finite.

To see that $\{p^{-1}(V_\lambda)\}$ is
a refinement of $\{U_\alpha\}$, note
that $p^{-1}(V_\lambda)\subset U_\alpha$
since $U_\alpha$ is invariant.
Moreover, $p^{-1}(V_\lambda)$ is
invariant.
\end{proof}

\noindent
{\it Note on the original.} Here $X$ is
a $G$-space and the proof uses that $G$
is compact, so that the orbit
projection is perfect; neither is in
the statement.

\begin{proposition}
\label{proposition-tubular-refinement}
Let $X$ be a paracompact $G$-space. If
$\{G(S_\alpha)\}_{\alpha\in A}$ is a
tubular cover (that is, $S_\alpha$ is
an $H_\alpha$-slice for $\alpha\in A$),
then there is a locally finite
refinement by tubular sets.
\end{proposition}

\begin{proof}
We know that $G(S_\alpha)$ is open in
$X$, so our cover is as in the previous
proposition. By the same construction
we obtain a refinement
$\{p^{-1}(V_\lambda)\}$. Take
$$
p^{-1}(V_\lambda)
\subset U_\alpha=G(S_\alpha)
$$
and define
$$
Q_\lambda
=p^{-1}(V_\lambda)\cap S_\alpha.
$$
In fact,
$p^{-1}(V_\lambda)=G(Q_\lambda)$, and
$Q_\lambda$ is an $H_\alpha$-slice.
\end{proof}

\begin{proposition}
\label{proposition-map-to-coset-space}
Suppose that $S$ is a global $H$-slice.
There is a function
$\tilde{g}:X\to G/H$ such that
$x\in gS$ for $g\in\tilde{g}(x)$, and
it is continuous, equivariant and
unique.
\end{proposition}

\begin{proof}
Since $X=G(S)$, any $x\in X$ is of the
form $gs$. Define $\tilde{g}(x)=gH$.

To see that $\tilde{g}$ is well
defined, take $g_x,h_x\in G$ such that
$x\in g_xS$ and $x\in h_xS$. Then
$g_x^{-1}x\in S$ and
$g_x^{-1}x\in g_x^{-1}h_xS$, and then
$$
S\cap g^{-1}_xh_xS\neq\varnothing.
$$
Since $S$ is a slice,
$g^{-1}_xh_x\in H$, that is,
$$
g_xH=h_xH=\tilde{g}(x).
$$

We see that $\tilde{g}:X\to G/H$ is
continuous. Let $x\in X$ and $U$ open
in $G/H$ be such that
$\tilde{g}(x)\in U$. Consider
$p:G\to G/H$, which is open, so
$p^{-1}(U)$ is open in $G$. Then
$W=E\cdot S$ is open in $X$. If
$x\in gS$, then $\tilde{g}(x)=gH$
\ldots\ {\bf (the rest is missing)}.
\end{proof}

\noindent
{\it Note on the original.} The proof
of continuity is unfinished there, and
the set $E$ is not defined; presumably
$E=p^{-1}(U)$, and $E\cdot S$ is then
the preimage of $U$ under $\tilde{g}$.
Its openness would follow from
Proposition
\ref{proposition-action-on-slice-open}
when $G$ is compact.

\begin{definition}
\label{definition-canonical-cover}
Let $U$ be an invariant open set of a
$G$-space $X$ and
$\{G(S_\lambda)\}_{\lambda\in\Lambda}$
a tubular cover, with $S_\lambda$ an
$H_\lambda$-slice. We say that this
cover is {\it $G$-canonical} if:
\begin{enumerate}
\item it is locally finite;
\item for each $\lambda\in\Lambda$
there is $x\in U$ such that
$G_x=H_\lambda$;
\item for each $a\in X\setminus U$ and
any neighborhood $V_a$ there is a
neighborhood $W_a$ of $a$ in $X$ such
that if $g\in G$ and
$gS_\lambda\cap W_a\neq\varnothing$
then
\begin{enumerate}
\item $gS_\lambda\subset V_a$,
\item there is $h\in G$ such that
$ha\in V_a$ and
$H_\lambda\subset g^{-1}G_{ha}g$.
\end{enumerate}
\end{enumerate}
\end{definition}

\noindent
{\it Note on the original.} The cover
is presumably a cover of $U$. In the
third condition the original writes
$gS_x$ for $gS_\lambda$ and
``$g^{-1}Ghag$'' for the conjugate of
the stabilizer; the reading
$g^{-1}G_{ha}g$ is mine.

\begin{proposition}[5.5 of Bredon, due
to Palais]
\label{proposition-bredon-palais}
Let $G$ be a compact Lie group. For
each neighborhood $U$ of $e$ in $G$ and
for every $x\in X$ there is an open set
$V$ of $X$ containing $x$ such that for
every $y\in V$ there is $u\in U$ such
that
$$
u^{-1}G_yu=G_x.
$$
\end{proposition}

\noindent
{\it Note on the original.} The
reference is given only as ``5.5 de
Bredon, de Palais''; I have not checked
it against the book. The conclusion
should be an inclusion,
$u^{-1}G_yu\subset G_x$, which is how
the result is used in the unfinished
proof below: stabilizers of nearby
points can be smaller.

\begin{theorem}
\label{theorem-canonical-cover-exists}
If $(X,d)$ is a metrizable $G$-space
with $G$ a compact Lie group, then
each invariant open set $U$ admits a
$G$-canonical cover.
\end{theorem}

\noindent
{\it Note on the original.} The source
file contains the following attempt at
a proof, disabled so that it does not
appear in the compiled notes. It is
unfinished, and it breaks down at the
place marked.

\medskip\noindent
{\it Unfinished proof.} If $U=X$, we
have that for every $x\in X$ there is a
$G_x$-slice. Then
$\{G(S_x)\}_{x\in X}$ is a tubular open
cover, for which, by
Proposition
\ref{proposition-tubular-refinement},
there is a refinement
$\{G(Q_\alpha)\}$ where $Q_\alpha$ is a
$G_x$-slice. Thus we have properties 1
and 2 of the definition. Property 3
holds vacuously.

Now suppose that $U\neq X$. Then $U$ is
metrizable, so it is paracompact by
Stone's theorem. Since $U$ is the union
of the orbits $G(x)$ with $x\in U$,
take a representative $x\in U$ of each
orbit and a $G_x$-slice $Q_x$. Define
$$
r_x=\frac{1}{4}
\rho(x,X\setminus U)
$$
which is a positive number because $U$
is open. Now define
$$
R_x=Q_x\cap B_{r_x}(x)
$$
which is a $G_x$-slice.

Then the $R_x$ {\it are an open cover}
of $U$, for which there is a tubular
refinement $\{U_\lambda\}$, with
$U_\lambda=G(S_\lambda)$ for
$S_\lambda$ a slice for
$G_x=H_\lambda$ of the form
$S_\lambda=R_x\cap U_x$. With this we
have proved the first two points of the
definition.

Let $a\in X\setminus U$ and $V_a$ open
in $X$ with $a\in V_a$. Since
$G\times X\to X$ is continuous and
$e\cdot a=a\in V_a$, there is $E$ open
in $G$ with $e\in E$ such that
$g\cdot a\in V_a$ for $g\in E$. Now we
apply
Proposition
\ref{proposition-bredon-palais}: there
is a ball $B_\varepsilon(a)$ such that
for every $y\in B_\varepsilon(a)$ there
is $u\in E$ with
$$
u^{-1}G_yu\subset G_a.
$$
Now let $W_a=B_{\varepsilon/4}(a)$ and
$g\in G$ be such that
$gS_\lambda\cap W_a\neq\varnothing$.
Recall that there is $x\in U$ such that
$S_x=R_x\cap U_x$ and $S_x$ is a
$G_x$-slice.

Now let
$y\in gS_x\cap W_a\subset gR_x$. Then
$$
\begin{aligned}
\rho(a,gx)
&\leq\rho(a,y)+\rho(y,gx) \\
&<\varepsilon/4+\rho(y,gx) \\
&\leq\varepsilon/4
 +\operatorname{diam}R_x
\end{aligned}
$$
and
$$
\begin{aligned}
\operatorname{diam}R_x
&\leq r_x \\
&=\frac{1}{2}\rho(x,X\setminus U) \\
&=\frac{1}{2}\rho(gx,X\setminus U) \\
&<\rho(gx,a)
\end{aligned}
$$
since $a\in X\setminus U$. Then
$$
\operatorname{diam}R_x
<\frac{1}{2}\rho(x,X\setminus U)
<\frac{1}{2}\rho(gx,a),
$$
hence
$$
\rho(a,gx)
<\varepsilon/4
+\frac{1}{2}\rho(gx,a),
$$
hence
$$
\rho(a,gx)<\varepsilon/2.
$$
Finally, {\bf here it failed}:
$$
\operatorname{diam}gR_x
\leq\frac{1}{2}
\rho(gx,X\setminus U)
\leq\frac{1}{2}
$$

\noindent
{\it Note on the original.} In this
attempt the original writes
$\rho(X,X\setminus U)$ and
$\rho(gX,X\setminus U)$, which I read
as distances from the points $x$ and
$gx$. The constant in $r_x$ changes
from $\frac{1}{4}$ to $\frac{1}{2}$,
the metric is used as if it were
invariant, and the ball has diameter at
most $2r_x$, not $r_x$.

\begin{theorem}[Dugundji]
\label{theorem-dugundji-canonical}
(The statement was left empty in the
original.)
\end{theorem}

\section{Proper actions of locally
compact groups}
\label{section-proper-actions}

\begin{definition}
\label{definition-proper-action}
\begin{enumerate}
\item Let $G$ be a locally compact
Hausdorff group and $X$ a $G$-space.
For $U,V\subset X$ we define the
{\it transporter} as
$$
\left\langle{} U,V \right\rangle{}
:=\{g\in G:gU\cap V\neq\varnothing\}.
$$
\item We say that
{\it $U$ is thin relative to $V$} if
the closure of
$\left\langle{} U,V \right\rangle{}$
is compact.
\item $U$ is called {\it small} if
every $x\in X$ has a neighborhood $V_x$
such that $V_x$ is thin relative to
$U$.
\item An action $G\curvearrowright X$
is called {\it proper} (in the sense of
Richard Palais) if every point
$x\in X$ has a small neighborhood
$V_x$. In this case we say that $X$ is
a {\it proper $G$-space}.
\end{enumerate}
\end{definition}

\begin{theorem}[van Dantzig]
\label{theorem-van-dantzig}
Let $X$ be a connected Riemannian
manifold. Then $\operatorname{Iso}(X)$
is a Lie group and the natural action
$$
\operatorname{Iso}(X)\times X\to X,
\qquad
(f,x)\mapsto f(x)
$$
is proper.
\end{theorem}

\noindent
{\it Note on the original.} The name is
written ``van Danzig'' there, with no
reference. That the isometry group of a
Riemannian manifold is a Lie group is
usually attributed to Myers and
Steenrod; I have not checked the
attribution of the statement as given.

\begin{remark}
\label{remark-van-dantzig-general}
The theorem can be generalized to a
locally compact, locally connected,
metrizable space and a locally compact
group.
\end{remark}

\begin{example}
\label{example-proper-actions}
\begin{enumerate}
\item $\mathbb{R}_+
\curvearrowright\mathbb{R}^2$ is not
proper because it has a fixed point.
\item Let $G$ be locally compact and
$H<G$. The action
$G\curvearrowright G/H$ is proper if
and only if $H$ is compact.
\item Let $X$ be a topological group
and $G<X$ locally compact. Then the
natural multiplicative action
$G\curvearrowright X$ is proper.
\end{enumerate}
\end{example}

\begin{theorem}[R. Palais, 1961]
\label{theorem-palais-invariant-metric}
Let $G$ be a Lie group that acts
properly on a separable metrizable
space $X$. Then there is a
$G$-invariant metric on $X$.
\end{theorem}

\begin{remark}
\label{remark-question-separability}
Question: is it possible to remove the
separability condition? This is an open
problem.
\end{remark}

\noindent
{\it Note on the original.} The problem
is called open in the notes; I have not
checked its current status.

The property of being $T_1$ passes to
the quotient $X\to X/G$ (the notes for
this are missing). What will happen
with the property of being $T_2$?

\begin{proposition}
\label{proposition-proper-hausdorff-quotient}
Let $G\curvearrowright X$ be proper. If
for any $x,y\in X$ there are relatively
thin neighborhoods $U_x$ and $U_y$,
then $X/G$ is $T_2$ if $X$ is $T_2$.
\end{proposition}

\begin{proof}
(Seen in class.)
\end{proof}

Using this proposition we prove that
the property $T_2$ also passes to the
quotient. Continuing with the
separation axioms, we have the
following proposition.

\begin{proposition}
\label{proposition-proper-tychonoff-quotient}
Let $G\curvearrowright X$ be proper. If
$X$ is $T_{3\frac{1}{2}}$, then $X/G$
is $T_{3\frac{1}{2}}$.
\end{proposition}

\bibliography{my}
\bibliographystyle{amsalpha}

\end{document}
